% =====================================================================
%  ERS/SRS -- SIMPA: Sistema Inteligente de Mantenimiento de Palma Africana
%  Entrega 4 (2B) -- Cierre del proyecto
%  Equipo AHMRV -- ISR-401 -- UTEQ -- Periodo 2026-2027 PPA
%
%  COMPILACION (ejecutar en este orden, desde 01_ERS/):
%     pdflatex ERS_SRS_2B_v2.0.tex
%     bibtex   ERS_SRS_2B_v2.0
%     pdflatex ERS_SRS_2B_v2.0.tex
%     pdflatex ERS_SRS_2B_v2.0.tex
%
%  NOTA: no se usa babel-spanish para garantizar la reproducibilidad
%  en cualquier instalacion base de TeX Live. Si el equipo tiene
%  texlive-lang-spanish puede descomentar la linea correspondiente.
% =====================================================================

\documentclass[11pt,a4paper]{article}

\usepackage[utf8]{inputenc}
\usepackage[T1]{fontenc}
% \usepackage[spanish,es-tabla]{babel}   % opcional
\renewcommand{\contentsname}{Tabla de contenido}
\renewcommand{\listfigurename}{Índice de figuras}
\renewcommand{\listtablename}{Índice de tablas}
\renewcommand{\tablename}{Tabla}
\renewcommand{\figurename}{Figura}
\renewcommand{\refname}{Referencias}

\usepackage{amsmath}   % requerido por \text{} en la fórmula del WSJF
\usepackage[margin=2.2cm]{geometry}
\usepackage{graphicx}
\usepackage{longtable}
\usepackage{array}
\usepackage{booktabs}
\usepackage{xcolor}
\usepackage{colortbl}
\usepackage{fancyhdr}
\usepackage{enumitem}
\usepackage{caption}
\usepackage{lastpage}
\usepackage{titlesec}
\usepackage{url}
\usepackage[hidelinks]{hyperref}

\definecolor{uteqverde}{RGB}{0,0,0}
\definecolor{grisclaro}{RGB}{242,242,242}

% \id: identificadores en monoespaciada con puntos de corte tras - _ y /
\newcommand{\idbreak}{\discretionary{}{}{}}
\newcommand{\id}[1]{\texttt{\begingroup\hyphenpenalty=10000\relax#1\endgroup}}
\renewcommand{\id}[1]{\texttt{#1}}
\newcolumntype{L}[1]{>{\raggedright\arraybackslash}p{#1}}
\newcolumntype{C}[1]{>{\centering\arraybackslash}p{#1}}

\sloppy
\setlength{\emergencystretch}{3em}
\setlength{\tabcolsep}{4pt}

% --- Control de figuras: ninguna excede el area util de la pagina
\usepackage[htt]{hyphenat}   % permite dividir identificadores en \texttt
\usepackage[export]{adjustbox}
\usepackage{float}
\usepackage{tikz}
\usetikzlibrary{shapes.geometric, arrows.meta, positioning, calc, fit}
\setkeys{Gin}{keepaspectratio,max width=\textwidth,max height=0.86\textheight}
% --- Mas tolerancia para que los floats no se acumulen al final
\renewcommand{\topfraction}{0.9}
\renewcommand{\bottomfraction}{0.8}
\renewcommand{\textfraction}{0.07}
\renewcommand{\floatpagefraction}{0.7}

\pagestyle{fancy}
\fancyhf{}
\lhead{\small ERS/SRS -- SIMPA v2.0}
\rhead{\small Equipo AHMRV -- ISR-401}
\cfoot{\small Página \thepage\ de \pageref{LastPage}}
\renewcommand{\headrulewidth}{0.4pt}

\titleformat{\section}{\normalfont\large\bfseries\color{uteqverde}}{\thesection.}{0.6em}{}
\titleformat{\subsection}{\normalfont\normalsize\bfseries}{\thesubsection}{0.6em}{}
\titleformat{\subsubsection}{\normalfont\small\bfseries}{\thesubsubsection}{0.6em}{}

% =====================================================================
%  MACRO PARA REQUISITOS FUNCIONALES (plantilla de 8 atributos)
%  Uso: \RF{ID}{Nombre}{Descripcion}{Actor/Origen}{Entradas/Salidas}
%          {Pre/Postcondiciones}{Prioridad}{Criterio de verificacion}
% =====================================================================
\newcommand{\RF}[8]{%
\begin{longtable}{|L{3.5cm}|L{12.0cm}|}
\hline
\rowcolor{grisclaro}\multicolumn{2}{|p{15.7cm}|}{\textbf{#1 --- #2}}\\
\hline
\textbf{Identificador} & \id{#1} \\ \hline
\textbf{Nombre} & #2 \\ \hline
\textbf{Descripción} & #3 \\ \hline
\textbf{Actor / Origen (evidencia)} & #4 \\ \hline
\textbf{Entradas / Salidas} & #5 \\ \hline
\textbf{Pre / Postcondiciones} & #6 \\ \hline
\textbf{Prioridad (MoSCoW)} & #7 \\ \hline
\textbf{Criterio de verificación} & #8 \\ \hline
\caption{#1 --- #2}
\end{longtable}
\vspace{-0.3cm}
}

\begin{document}

% =====================================================================
%  PORTADA
% =====================================================================
\begin{titlepage}
\centering
\vspace*{0.8cm}
{\Large\bfseries UNIVERSIDAD TÉCNICA ESTATAL DE QUEVEDO\par}
\vspace{0.3cm}
{\large Facultad de Ciencias de la Computación\par}
{\large Carrera de Software (Rediseño)\par}
\vspace{1.2cm}
\rule{\textwidth}{1.2pt}\par
\vspace{0.5cm}
{\LARGE\bfseries\color{uteqverde} ESPECIFICACIÓN DE REQUISITOS DE SOFTWARE\par}
\vspace{0.4cm}
{\Large\bfseries SIMPA --- Sistema Inteligente de Mantenimiento\\[0.2cm] de Palma Africana\par}
\vspace{0.4cm}
{\large Entrega 4 (2B) --- Cierre del proyecto\par}
\vspace{0.5cm}
\rule{\textwidth}{1.2pt}\par
\vspace{1.0cm}

{\large\bfseries Organización cliente:\par}
{\large Palmicultora M (seudónimo)\par}
{\normalsize Cantón El Empalme, provincia del Guayas, Ecuador\par}
\vspace{0.8cm}

{\bfseries Equipo AHMRV --- Paralelo 4to ``A''\par}
\vspace{0.3cm}
\begin{tabular}{ll}
Villafuerte Rosero Allan Noe & Analista Líder \\
Huilcapi León Denisses Fabiola & Documentadora \\
Rizzo Vélez Edson Nagib & Verificador \\
Macías Herrera Josthyn Esteban & Modelador \\
Arboleda Yanza Francisco Javier & Apoyo Modelado \\
Alcívar Vélez Anderson Adonis & Apoyo Repositorio \\
\end{tabular}

\vspace{0.9cm}
{\bfseries Asignatura:} Ingeniería de Requerimientos (ISR-401) --- 4to Nivel\par
{\bfseries Docente:} Ing. Gleiston Cicerón Guerrero Ulloa, PhD\par
{\bfseries Estándar:} ISO/IEC/IEEE 29148:2018\par
{\bfseries Versión:} 2.0 \qquad {\bfseries Fecha:} 31/08/2026

\vspace{0.8cm}
{\bfseries Repositorio del proyecto:}\par
\url{https://github.com/gleiston-guerrero/SIMPA_ISR401}

\vfill
{\small Quevedo --- Los Ríos --- Ecuador\par}
\end{titlepage}

% =====================================================================
%  HISTORIAL DE VERSIONES
% =====================================================================
\section*{Historial de versiones}
\addcontentsline{toc}{section}{Historial de versiones}

\begin{longtable}{|C{1.4cm}|C{2.2cm}|L{3.0cm}|L{7.6cm}|}
\hline
\rowcolor{grisclaro}\textbf{Revisión interna} & \textbf{Fecha} & \textbf{Autor} & \textbf{Descripción del cambio} \\
\hline
\endhead
1.0 & 01/06/2026 & Equipo AHMRV & Elaboración inicial de la Entrega 1 (1A): planificación y elicitación. \\ \hline
1.1 & 23/06/2026 & Equipo AHMRV & Incorporación de correcciones docentes de la 1A: diagrama de contexto, actas formales, trazabilidad de requisitos brutos. \\ \hline
2.0 & 26/06/2026 & Equipo AHMRV & Entrega 2 (1B): requisitos específicos, modelado UML, priorización MoSCoW y matriz de trazabilidad parcial. \\ \hline
3.0 \newline \textbf{(= ERS v1.0)} & 03/08/2026 & Equipo AHMRV & \textbf{Entrega 3 (2A).} Correcciones del informe docente de la 1B incorporadas (véase §1.6). Segunda ronda de campo con cinco nuevos participantes (EV-04 a EV-08) y una nueva organización fuente. Ampliación a 39 RF y 18 RNF sobre las nueve características de ISO/IEC 25010:2023. Incorporación del requisito de explicabilidad, del mapeo a la LOPDP, de las historias de usuario en formato Connextra con escenarios Gherkin y del modelado organizacional i*. Renumeración del catálogo de evidencias (véase §1.7). Migración del documento a \LaTeX. \\ \hline
\rowcolor{grisclaro}3.1 \newline \textbf{(= ERS v1.1)} & 05/08/2026, actualizada el 17/08/2026. & Equipo AHMRV & \textbf{Línea base aprobada por el Change Control Board (PE4, Unidad IV).} Corrección del 100\% de los defectos de severidad crítica y mayor detectados en la inspección formal tipo Fagan de la §1 a la §7 (registro completo en el Anexo B del informe PE4). Se consolida la incorporación de la evidencia \id{EV-13}, cuya recolección concluyó el 07/08/2026, es decir, con posterioridad al cierre de la versión 3.0 el 03/08/2026: los requisitos \id{RF-36} a \id{RF-39} se habían incorporado a la v3.0 sobre una evidencia todavía en curso, y es esta versión la que los sustenta sobre la evidencia completa. Se implementan además los cambios aprobados en las solicitudes \id{RFC-01}, \id{RFC-02} y \id{RFC-03}.
\textbf{Actualización de cierre (PE5, Unidad V):} auditoría de calidad con las seis métricas derivadas de
ISO/IEC 25010:2023 y corrección de toda métrica por debajo de su valor de referencia; re-inspección sobre
el documento corregido, con registro de los defectos residuales y su cierre; especificación de los
componentes de inteligencia artificial del sistema con sus requisitos de rendimiento, equidad y
explicabilidad; incorporación de los diagramas de flujo de datos y cierre de la trazabilidad extremo a
extremo. En PE5, esta revisión quedó declarada como la línea base ERS v1.1 del proyecto integrador. \\ \hline
3.2 \newline \textbf{(= ERS v2.0)} & 31/08/2026 & Equipo AHMRV & \textbf{Entrega 4 (2B) --- cierre del proyecto.} Migración y consolidación de PE5 en el repositorio principal: catálogo vigente de 42 RF y 19 RNF, 24 historias de usuario, casos de uso CU-01 a CU-18, matriz extremo a extremo de 73 filas, priorización post-CCB de 24 Must, 16 Should y 2 Could, corrección de la traza de RNF-11 a EV-07 y EV-10, e integración de los anexos de auditoría, CCB y declaración de uso de IA. \\ \hline
\caption{Historial de versiones del documento}
\end{longtable}

\vspace{0.4cm}
\noindent\fbox{\parbox{0.97\textwidth}{%
\small\textbf{Dos numeraciones, una línea base vigente.} La columna anterior registra las
\emph{revisiones internas} con que el equipo ha venido versionando el documento desde la Entrega 1A.
La asignatura emplea una numeración distinta: designó como \textbf{ERS v1.0} al documento de la
Entrega 2A ---revisión interna 3.0--- y como \textbf{ERS v1.1} a la versión resultante del Change
Control Board de la Práctica Experimental de la Unidad IV ---revisión interna 3.1---. Tras la
consolidación de la Entrega 4 (2B), la versión oficial vigente es \textbf{ERS v2.0}. Esta línea base
se identifica en este repositorio mediante la etiqueta anotada \id{baseline-v2.0}. Las revisiones
anteriores se conservan porque documentan la evolución real del documento y no deben interpretarse
como la versión vigente.}}

\vspace{0.3cm}
\noindent\fbox{\parbox{0.97\textwidth}{%
\small\textbf{Regla de acumulación.} Este documento contiene y reemplaza a las Entregas 1 (1A) y 2 (1B).
Las secciones heredadas fueron revisadas e incorporan las observaciones del docente registradas en el
informe individual de la Entrega 2 (1B), fechado el 01/07/2026. La trazabilidad de dichas correcciones
se detalla en la sección §1.6.}}

\newpage
\tableofcontents
\newpage
\listoffigures
\listoftables
\newpage

% =====================================================================
%  1. INTRODUCCION
% =====================================================================
\section{Introducción}

\subsection{Propósito del documento}

El presente documento especifica los requisitos funcionales, no funcionales y legales del
\textbf{Sistema Inteligente de Mantenimiento de Palma Africana (SIMPA)}, siguiendo la estructura
del estándar ISO/IEC/IEEE 29148:2018 \cite{iso29148}. Su propósito es establecer una comprensión
común, verificable y sin ambigüedad de lo que el sistema debe hacer, de las restricciones bajo las
que debe operar y de la evidencia de campo que sustenta cada decisión de especificación.

El documento está dirigido a: el \textbf{cliente} (propietario y administración de la organización
Palmicultora M), el \textbf{equipo de desarrollo} responsable del diseño e implementación, los
\textbf{verificadores} encargados de validar el cumplimiento de cada criterio de aceptación, y los
\textbf{revisores académicos} de la asignatura ISR-401.

La redacción de los requisitos sigue los criterios de calidad de Pohl \cite{pohl2010} y de Wiegers y
Beatty \cite{wiegers2013}, y el vocabulario del cuerpo de conocimiento SWEBOK v4.0 \cite{swebok4} y
del sílabo CPRE \cite{ireb2023}. La atención a la verificabilidad y a la ausencia de ambigüedad
responde a la evidencia empírica de que estos son los defectos de requisitos más costosos en la
práctica industrial \cite{mendez2017,montgomery2022}, y de que la ambigüedad se origina con
frecuencia en el conocimiento tácito no explicitado durante las entrevistas de elicitación
\cite{ferrari2016}.

A diferencia de las entregas anteriores, esta versión incorpora el principio de \emph{trazabilidad
evidencial estricta}: ningún requisito de este documento se enuncia sin la referencia explícita al
identificador de evidencia (\id{EV-XX}) que lo respalda, y toda evidencia referenciada reposa
físicamente dentro del repositorio del proyecto.

\subsection{Alcance del producto}

\subsubsection{Nombre y descripción general}

\textbf{SIMPA} es una aplicación móvil Android con respaldo en la nube, destinada a la gestión,
el monitoreo y el diagnóstico asistido del cultivo de palma africana (\emph{Elaeis guineensis} y
sus híbridos interespecíficos) en una explotación agrícola de aproximadamente cien hectáreas.

El sistema centraliza información que actualmente se registra de forma manual en libretas y se
transmite de forma informal por mensajería instantánea, e incorpora componentes de inteligencia
artificial para el análisis de imágenes de la palma y de sus racimos.

\subsubsection{Alcance de esta versión}

SIMPA v1 opera sobre \textbf{una sola unidad productiva}: la finca Palmicultora M. La organización
cliente administra además una segunda finca bajo la misma dirección técnica, cuyos documentos
operativos se emplean en este documento como evidencia del proceso de planificación
(\id{EV-11}); no obstante, la operación multi-finca se declara explícitamente como
\textbf{trabajo futuro} y no forma parte del alcance de esta versión.

\subsubsection{Funcionalidades principales}

\begin{itemize}[nosep]
  \item Gestión de plantaciones, lotes y variedades de palma.
  \item Planificación semanal de labores con seguimiento de cumplimiento y arrastre de pendientes.
  \item Registro de labores agrícolas por unidad de avance (coronas, bancos, plantas, jornales).
  \item Registro del proceso de polinización asistida y conteo georreferenciado de flores polinizadas.
  \item Análisis de imágenes mediante IA para detección de plagas y enfermedades, diagnóstico de
        deficiencias nutricionales y clasificación de madurez del racimo.
  \item Generación de alertas tempranas con explicación comprensible para personal no especializado.
  \item Registro de la calificación de calidad emitida por la planta extractora y trazabilidad
        del lote rechazado hasta la cuadrilla responsable.
  \item Estimación de producción y generación de reportes diarios, semanales y anuales.
  \item Ciclo de liquidación semanal: catálogo de tarifas por labor, cálculo de la remuneración a
        partir del avance registrado, comprobante individual y consolidado exportable (\id{RF-36} a
        \id{RF-39}). La §7.2 lo identifica como proceso central de la organización.
\end{itemize}

\subsubsection{Exclusiones explícitas}

El sistema \textbf{no} realizará: control automático del sistema de riego; aplicación automática de
fertilizantes o pesticidas; operación de drones o maquinaria agrícola; gestión de nómina, contabilidad
o facturación; comercialización de la producción; ni el procesamiento industrial del fruto en la
planta extractora.

\subsection{Definiciones, acrónimos y abreviaturas}

\begin{longtable}{|L{3.6cm}|L{11.8cm}|}
\hline
\rowcolor{grisclaro}\textbf{Término} & \textbf{Definición} \\
\hline
\endhead
Antesis & Etapa de floración en la que la flor femenina está receptiva para la polinización; se reconoce por coloración amarilla o negra (\id{EV-03}). \\ \hline
Banco & Unidad de superficie despejada alrededor de la base de la palma; unidad de medida de avance en labores de limpieza (\id{EV-11}). \\ \hline
Chapia & Labor de corte de maleza en la calle del cultivo. \\ \hline
Corona & Área circular limpia alrededor del tronco de la palma; también, la labor de mantenerla despejada. \\ \hline
Coarí × La Mé & Híbrido interespecífico de palma aceitera, resistente a la pudrición del cogollo; variedad sembrada en Palmicultora M (\id{EV-02}). \\ \hline
Desprendimiento & Caída espontánea de frutos del racimo; indicador operativo de madurez usado por la industria extractora (\id{EV-04}, \id{EV-08}). \\ \hline
Estrategus & Coleóptero que perfora el suelo y daña el meristema de la palma joven (\id{EV-02}). \\ \hline
Explicabilidad & Capacidad de un sistema con IA de comunicar a la persona usuaria las razones de una decisión automática, en términos que dicha persona pueda comprender \cite{chazette2020}. \\ \hline
Flecha & Hoja no abierta en el ápice de la palma; su acumulación indica déficit hídrico (\id{EV-01}). \\ \hline
Guía de remisión & Documento de transporte exigido por la normativa ecuatoriana para el traslado de mercancía; su ausencia expone al productor a sanción (\id{EV-04}). \\ \hline
Guineensis & \emph{Elaeis guineensis}, variedad africana de palma aceitera; umbral de madurez de 3 frutos desprendidos (\id{EV-04}). \\ \hline
Jornal & Unidad de trabajo equivalente a una jornada-persona; unidad de planificación en los documentos de la organización (\id{EV-11}). \\ \hline
Lote & Subdivisión operativa de la plantación; unidad de asignación de labores y de trazabilidad. \\ \hline
LOPDP & Ley Orgánica de Protección de Datos Personales del Ecuador \cite{lopdp}. \\ \hline
Meristema & Tejido de crecimiento de la palma; su daño provoca la muerte de la planta. \\ \hline
MoSCoW & Técnica de priorización: \emph{Must}, \emph{Should}, \emph{Could}, \emph{Won't}. \\ \hline
Pedúnculo & Base leñosa del racimo; si supera 5 cm absorbe aceite durante la extracción (\id{EV-04}). \\ \hline
Racimo (RFF) & Racimo de fruta fresca; unidad de cosecha y de comercialización. \\ \hline
RNF & Requisito no funcional; criterio de calidad según ISO/IEC 25010:2023 \cite{iso25010}. \\ \hline
Ticket de pesaje & Comprobante emitido por la extractora con peso bruto, tara, neto y calificación de calidad del lote entregado (\id{EV-04}). \\ \hline
\caption{Glosario del dominio del problema}
\end{longtable}

\subsection{Referencias}

Las referencias bibliográficas se listan al final del documento en formato IEEE y se gestionan
íntegramente en el archivo \id{referencias.bib}. El documento cita 27 fuentes primarias, de las
cuales 12 corresponden al período 2023--2026, conforme al mínimo exigido de 25 fuentes con al
menos 10 recientes.

\subsection{Nota sobre el uso de herramientas de asistencia en la redacción}

Conforme a la advertencia global de la guía entonces aplicable a la Entrega 3 (2A), se deja constancia de que toda
afirmación sustantiva de este documento está anclada a evidencia primaria del repositorio o a
referencia bibliográfica revisada por pares. El uso de modelos grandes de lenguaje se limitó a
tareas de asistencia en la redacción y el formato. Esta delimitación es coherente con las guías
sistemáticas de uso de estos modelos en tareas de ingeniería de requisitos
\cite{vogelsang2024,arora2024} y con la revisión sistemática de Cheng \emph{et al.}
\cite{cheng2026}, que constituye además el marco del componente empírico descrito en la §6.

\subsection{Visión general del documento}

La organización del documento sigue la estructura de ISO/IEC/IEEE 29148:2018 \cite{iso29148},
complementada con las prácticas de especificación descritas por Sommerville \cite{sommerville2016}.
La sección §2 describe el contexto del sistema, sus usuarios y el modelado organizacional. La §3
especifica los requisitos funcionales, no funcionales, legales y las restricciones de diseño. La §4
presenta el modelado UML. La §5 desarrolla la priorización y la trazabilidad extendida. La §6
describe el prototipo funcional. Los apéndices contienen el plan de ingeniería de requisitos, las
evidencias de campo, la clasificación completa y la matriz de trazabilidad.

\subsection{Incorporación de las observaciones de la Entrega 2 (1B)}

Conforme al gatekeeper G2 de la guía entonces aplicable a la Entrega 3 (2A), este documento parte de la 1B corregida. La
siguiente tabla traza cada observación del informe docente del 01/07/2026 con la acción adoptada.

\begin{longtable}{|C{1.0cm}|L{6.6cm}|L{7.8cm}|}
\hline
\rowcolor{grisclaro}\textbf{\#} & \textbf{Observación del docente (1B)} & \textbf{Acción adoptada en la 2A} \\
\hline
\endhead
1 & Audio de la entrevista al administrador general declarado solo mediante enlace externo a YouTube, sin archivo dentro del repositorio. & Archivo de audio incorporado físicamente al repositorio mediante Git LFS; ficha técnica con duración, códec y hash SHA-256 publicada en zona pública. \\ \hline
2 & Videos de dos entrevistas declarados solo mediante enlaces externos. & Videos incorporados como archivos dentro del contenedor cifrado \id{02\_Evidencias/00\_Restringido/}, con hash registrado en \id{checksums.sha256}. \\ \hline
3 & Ausencia de fotografías de la aplicación real del cuestionario. & \id{02\_Evidencias/Cuestionario/Fotos\_Aplicacion/} contiene 15 capturas de pantalla del instrumento (commits \id{34ce578} y \id{0792213}) y 5 fotografías de la sesión de consentimiento complementario del 15/09/2026, no fotografías de la aplicación real en campo con participantes fuera de encuadre. La carpeta y su README declaran esta corrección (tarea H2). \\ \hline
4 & Ausencia de documentos originales de la organización. & Incorporados los planes semanales de labores de dos unidades productivas (\id{EV-11}), anonimizados conforme a la §4.1 de la guía. \\ \hline
5 & Archivo de evidencia fotográfica con tamaño de 2 bytes (\emph{placeholder}). & Archivo reemplazado por la imagen real y verificado mediante inspección de cabecera binaria. \\ \hline
6 & Consentimientos publicados con cédula y firma visibles. & Copias públicas enmascaradas conforme a la §4.1; originales íntegros trasladados a la zona restringida cifrada. \\ \hline
\caption{Trazabilidad de las correcciones de la Entrega 2 (1B)}
\end{longtable}

\subsection{Nota sobre la renumeración del catálogo de evidencias}

La segunda ronda de campo incorporó cinco nuevas entrevistas. Para mantener la continuidad de las
entrevistas de la primera ronda (\id{EV-01} a \id{EV-03}), las evidencias no-entrevista de la 1B se
desplazaron según la siguiente equivalencia, registrada también en \id{CHANGELOG.md}:

\begin{center}
\begin{tabular}{|c|c|l|}
\hline
\rowcolor{grisclaro}\textbf{ID en la 1B} & \textbf{ID en la 2A} & \textbf{Contenido} \\ \hline
\id{EV-04} & \id{EV-09} & Observación de campo \\ \hline
\id{EV-05} & \id{EV-10} & Cuestionario, aplicación piloto ($n=4$) \\ \hline
\end{tabular}
\end{center}

\vspace{0.3cm}
\noindent Las referencias \id{EV-04} a \id{EV-08} corresponden ahora a las entrevistas de la segunda
ronda. El catálogo completo se detalla en el Apéndice B.1.

% =====================================================================
%  2. DESCRIPCION GENERAL
% =====================================================================
\newpage
\section{Descripción general}

\subsection{Perspectiva del producto y límite del sistema}

SIMPA es un sistema \textbf{nuevo}. No reemplaza a un sistema informático previo: sustituye un
conjunto de prácticas manuales. La evidencia de campo documenta el proceso actual (AS-IS) de la
siguiente manera:

\begin{itemize}[nosep]
  \item El registro diario de avance se anota en una libreta personal por cada trabajador y se
        entrega verbalmente al encargado al final de la jornada (\id{EV-05}, \id{EV-06}).
  \item El reporte hacia la administración se transmite por mensajería instantánea; la
        organización no emplea ninguna otra herramienta digital (\id{EV-07}).
  \item La detección de anomalías depende de la inspección visual y de la experiencia del personal
        técnico; el recorrido completo de la plantación toma una semana (\id{EV-01}).
  \item La supervisión presencial del asistente de administración ocurre dos veces por semana, lo
        que introduce una latencia de hasta tres días entre el aviso y la verificación (\id{EV-07}).
\end{itemize}

El límite del sistema se representa en la Figura~\ref{fig:contexto}. SIMPA intercambia información
con cinco entidades externas: el dispositivo móvil (cámara y GPS), la estación meteorológica, la
base de conocimiento de plagas, el servicio de análisis de imágenes y la \textbf{planta extractora},
actor incorporado en esta entrega a partir de \id{EV-04} y \id{EV-08}.

\begin{figure}[!htbp]
\centering
\includegraphics[width=1.0\textwidth]{img/contexto}
\caption{Diagrama de contexto del SIMPA (límite del sistema).}
\label{fig:contexto}
\end{figure}

\subsection{Funciones del producto}

Las funciones de alto nivel se agrupan en siete bloques: (i) gestión de la estructura productiva;
(ii) planificación y seguimiento de labores; (iii) registro operativo de campo; (iv) diagnóstico
asistido por IA; (v) control de calidad y relación con la extractora; y (vi) reportes y estimación
de producción. El detalle se especifica en la §3.2.

\subsection{Clases y características de los usuarios}

\begin{longtable}{|L{3.0cm}|L{5.4cm}|L{3.2cm}|L{2.8cm}|}
\hline
\rowcolor{grisclaro}\textbf{Perfil} & \textbf{Descripción} & \textbf{Competencia digital} & \textbf{Frecuencia de uso} \\
\hline
\endhead
Administrador general & Supervisión global del cultivo, decisiones de inversión y evaluación del personal (\id{EV-01}). & Media & Diaria \\ \hline
Administrador / Asesor técnico & Dirección técnica del cultivo; define nutrición, sanidad y manejo agrícola (\id{EV-02}). & Media & Semanal \\ \hline
Asistente de administración & Seguimiento de los trabajos de campo, verificación fotográfica de afectaciones, visitas dos veces por semana (\id{EV-07}). & \textbf{Alta} & Bisemanal \\ \hline
Jefe de polinización & Dirige la polinización asistida y controla su calidad (\id{EV-03}). & Media & Diaria \\ \hline
Trabajador agrícola & Ejecuta chapia, corona, polinización y cosecha; registra su avance (\id{EV-05}, \id{EV-06}). & \textbf{Baja} & Diaria \\ \hline
Técnico de la extractora & Califica la fruta recibida y asesora al productor sobre calidad (\id{EV-04}, \id{EV-08}). & Media-alta & Semanal \\ \hline
\caption{Clases de usuarios del SIMPA}
\end{longtable}

\vspace{0.2cm}
\noindent\fbox{\parbox{0.97\textwidth}{\small\textbf{Clase de usuario y parte interesada no son lo mismo.}
Esta tabla enumera las \emph{seis clases de usuario} del sistema: los perfiles que operan la
aplicación y que, por tanto, son actor de al menos un caso de uso. La §2.4 enumera además las
\emph{partes interesadas}, conjunto más amplio que incluye a quienes tienen poder o interés sobre el
proyecto sin interactuar con el sistema. Dos de ellas ---el propietario de la organización y la
planta extractora como entidad--- no son clase de usuario: el primero recibe los reportes a través
de la administración general y la segunda interactúa mediante su técnico, que sí es clase de
usuario. La distinción se declara aquí de forma expresa porque la auditoría de calidad de la
Unidad~V mide la cobertura de actores sobre las clases de usuario, no sobre las partes interesadas,
y confundir ambos conjuntos produce una métrica falsamente incumplida.}}

\noindent
La brecha de competencia digital entre el asistente de administración y el trabajador agrícola es
un hallazgo relevante para el diseño. El primero declaró estar familiarizado con este tipo de
tecnología y anticipó que la dificultad recaería sobre el personal de campo: \emph{``Más sería el equipo de campo
que hace las labores manuales como fertilización, polinización y poda.''} (\id{EV-07}). Esta observación sustenta la restricción
\id{RD-06} y el requisito de usabilidad \id{RNF-08}.

\subsection{Mapa de partes interesadas (matriz poder/interés)}

\begin{longtable}{|L{4.2cm}|C{1.8cm}|C{1.8cm}|L{3.4cm}|C{2.2cm}|}
\hline
\rowcolor{grisclaro}\textbf{Parte interesada} & \textbf{Poder} & \textbf{Interés} & \textbf{Estrategia} & \textbf{Evidencia} \\
\hline
\endhead
Propietario de la organización & Alto & Alto & Gestionar de cerca & \id{EV-02} \\ \hline
Administrador general & Alto & Alto & Gestionar de cerca & \id{EV-01} \\ \hline
Administrador / Asesor técnico & Alto & Alto & Gestionar de cerca & \id{EV-02} \\ \hline
Asistente de administración & Medio & Alto & Mantener informado & \id{EV-07} \\ \hline
Jefe de polinización & Medio & Alto & Mantener informado & \footnotesize \id{EV-03} \\ \hline
Planta extractora & \textbf{Alto} & Medio & \textbf{Mantener satisfecho} & \id{EV-04}, \id{EV-08} \\ \hline
Trabajadores agrícolas & Bajo & Alto & Mantener informados & \id{EV-05}, \id{EV-06} \\ \hline
Agrocalidad (ente regulador) & Alto & Bajo & Monitorear & \footnotesize \id{EV-04} \\ \hline
\caption{Mapa de partes interesadas con clasificación poder/interés}
\end{longtable}

\noindent
La incorporación de la planta extractora como parte interesada de \textbf{poder alto} es la
modificación más significativa respecto de la 1B. La extractora fija unilateralmente los criterios
de aceptación de la fruta, aplica castigos económicos sobre el lote y puede rechazar la carga
completa (\id{EV-04}). Su clasificación como parte de alto poder e interés medio refleja que, si
bien no es usuaria directa del SIMPA, sus criterios determinan el valor económico de la producción.

\subsection{Modelado organizacional i*}

\subsubsection{Diagrama de Dependencia Estratégica (SD)}

El modelo SD representa las dependencias entre actores del dominio, independientemente del sistema.
La Tabla~\ref{tab:sd} resume las dependencias identificadas a partir de la evidencia.

\begin{longtable}{|L{3.4cm}|L{3.4cm}|L{2.4cm}|L{4.4cm}|}
\hline
\rowcolor{grisclaro}\textbf{Actor dependiente} & \textbf{Actor depositario} & \textbf{Tipo} & \textbf{Elemento de dependencia} \\
\hline
\endhead
Propietario & Administrador general & Objetivo & Rendimiento de 40--45 t/ha/año (\id{EV-02}) \\ \hline
Administrador general & Jefe de campo & Recurso & Reporte diario de avance (\id{EV-07}) \\ \hline
Administrador general & Trabajador agrícola & Tarea & Ejecución de la labor presupuestada (\id{EV-11}) \\ \hline
Administrador general & Asistente de administración & Objetivo-blando & Oportunidad en la detección de afectaciones (\id{EV-07}) \\ \hline
Trabajador agrícola & Administrador general & Recurso & Asignación de lote y pago por avance (\id{EV-05}, \id{EV-06}) \\ \hline
Extractora & Cuadrilla de cosecha & Objetivo-blando & Fruta con desprendimiento suficiente (\id{EV-04}) \\ \hline
Extractora & Productor & Recurso & Estimación de tonelaje para planificar capacidad (\id{EV-04}) \\ \hline
Productor & Extractora & Recurso & Ticket de pesaje con calificación de calidad (\id{EV-04}) \\ \hline
Jefe de polinización & Polinizador & Tarea & Aplicación correcta del polinizante (\id{EV-03}) \\ \hline
\caption{Dependencias estratégicas (modelo i* SD)}
\label{tab:sd}
\end{longtable}

\noindent\textbf{Nota sobre la notación.} Los modelos i* se representan con la notación estándar de
Yu: objetivo (\emph{goal}) en rectángulo redondeado, objetivo-blando (\emph{softgoal}) en nube,
tarea (\emph{task}) en hexágono y recurso (\emph{resource}) en rectángulo. En el modelo SD, la
lectura de cada dependencia es \emph{depender} $\rightarrow$ \emph{dependum} $\rightarrow$
\emph{dependee}. Cada diagrama incorpora su leyenda.

\begin{figure}[!htbp]
\centering
\includegraphics[width=1.0\textwidth]{img/istar_sd}
\caption{Diagrama de Dependencia Estratégica (i* SD) del dominio de Palmicultora M.}
\label{fig:sd}
\end{figure}

\subsubsection{Diagrama de Razón Estratégica (SR)}

El modelo SR profundiza en el actor \textbf{Administrador general}, cuyo objetivo raíz
\emph{``mantener el cultivo en estado productivo''} se descompone en cuatro tareas: supervisar
equipos de trabajo, controlar el estado sanitario, asegurar la humedad del suelo y verificar el
cumplimiento del presupuesto semanal. Sobre estas tareas actúan tres objetivos-blandos derivados de
la evidencia: \emph{oportunidad del diagnóstico} (\id{EV-01}: el recorrido completo toma una
semana), \emph{confiabilidad del conteo} (\id{EV-03}: el conteo manual afecta el rendimiento del
polinizador) y \emph{reducción de la dependencia del conocimiento experto} (\id{EV-01},
\id{EV-02}: la escasez de mano de obra calificada es el principal problema de gestión).

\begin{figure}[!htbp]
\centering
\includegraphics[width=1.0\textwidth]{img/istar_sr}
\caption{Diagrama de Razón Estratégica (i* SR) del actor Administrador general.}
\label{fig:sr}
\end{figure}

\subsection{Entorno operativo}

\textbf{Entorno físico.} Plantación de aproximadamente cien hectáreas dividida en seis lotes
(\id{EV-11}), con campamento y oficina administrativa. Condiciones de exposición solar directa,
humedad alta y jornadas que se extienden entre las 06:00 y las 18:00 según el clima (\id{EV-06}).

\textbf{Conectividad.} Existe internet satelital parcial; la conexión estable se limita al
campamento (\id{EV-07}). La aplicación ampliada del cuestionario ($n=62$, \id{EV-12}) precisa esta
condición: el 25{,}8\% de las personas respondientes declara disponer de señal siempre, el 51{,}6\%
solo a veces, el 21{,}0\% casi nunca y el 1{,}6\% nunca. Es decir, \textbf{el 74{,}2\% no cuenta con
conectividad permanente} en su zona de trabajo. Esta distribución --- y no la ausencia total de
señal --- es la que determina la arquitectura \emph{offline-first} del sistema: la aplicación debe
ser plenamente funcional sin conexión y sincronizar de forma oportunista.

\textbf{Plataforma.} Aplicación móvil Android para campo; panel web para administración; backend en
la nube con API REST y servicio de inferencia de IA.

\subsection{Suposiciones y dependencias}

\textbf{Suposiciones.} Las imágenes capturadas tienen calidad suficiente para el análisis; la
organización mantiene la práctica de planificación semanal documentada (\id{EV-11}); la
organización proveerá un mecanismo de registro asistido para el personal sin dispositivo propio.

\vspace{0.2cm}
\noindent\fbox{\parbox{0.97\textwidth}{%
\small\textbf{Suposición invalidada por la evidencia.} Las entregas anteriores asumían que la
totalidad del personal disponía de teléfono inteligente propio. La aplicación ampliada del
cuestionario ($n=62$, \id{EV-12}) refuta esa suposición: \textbf{el 11{,}3\% de las personas
respondientes declara no usar teléfono inteligente}. Esta constatación motiva la restricción
\id{RD-10} y el requisito \id{RF-35}, que habilitan el registro delegado para dicho segmento. De no
haberse corregido, el sistema habría excluido a aproximadamente una de cada nueve personas
destinatarias.}}

\textbf{Dependencias.} Disponibilidad del servicio de inferencia de IA; disponibilidad de datos
climáticos; vigencia del consentimiento informado de los trabajadores para el tratamiento de datos
de geolocalización.

\textbf{Riesgos.} Baja calidad de imagen en condiciones de contraluz; resistencia al uso de
tecnología por parte del personal de campo (\id{EV-07}); variabilidad del criterio de madurez entre
variedades, que exige parametrización (\id{EV-04}).


\subsection[Resultados del cuestionario aplicado (n=62)]{Resultados del cuestionario aplicado ($n=62$)}

El cuestionario se aplicó al personal de campo en dos rondas con \textbf{instrumento idéntico}, con
el fin de preservar la comparabilidad: una aplicación piloto de cuatro respondientes
(\id{EV-10}, 30/06/2026) y una aplicación ampliada de 62 respondientes (\id{EV-12}).

\textbf{Definición del perfil dominante.} Para la Entrega 4 (2B), la guía vigente establece
$n \geq 60$ por perfil dominante, o un cálculo de potencia que justifique un tamaño menor.
El cuestionario reúne $n = 62$ respuestas. A nivel ocupacional amplio,
\textbf{60 participantes} corresponden a labores agrícolas de campo y 2 a la categoría
``Otra''; por tanto, el perfil agregado de personal agrícola de campo alcanza
$n = 60$ (96{,}8\% de la muestra) bajo esta definición.

Dentro de ese perfil, los subperfiles según la labor principal son polinización
($n = 18$), control fitosanitario ($n = 18$), labores de mantenimiento
($n = 13$) y riego ($n = 11$). Ningún subperfil individual alcanza por separado
el umbral de 60. En consecuencia, las inferencias se limitan al perfil ocupacional
agregado y cualquier análisis comparativo por labor tiene carácter exploratorio.
No se declara cumplimiento del umbral a nivel de subperfil.

\begin{longtable}{|L{6.8cm}|C{2.0cm}|C{2.0cm}|L{3.6cm}|}
\hline
\rowcolor{grisclaro}\textbf{Variable} & \textbf{Valor} & \textbf{\%} & \textbf{Implicación} \\
\hline
\endhead
\multicolumn{4}{|l|}{\textbf{Perfil de la muestra}} \\ \hline
Polinización & 18 & 29{,}0 & Estrato mayor \\ \hline
Control fitosanitario & 18 & 29{,}0 & Estrato mayor \\ \hline
Labores de mantenimiento & 13 & 21{,}0 & --- \\ \hline
Riego & 11 & 17{,}7 & --- \\ \hline
Otra & 2 & 3{,}2 & Excluido del perfil dominante \\ \hline
\multicolumn{4}{|l|}{\textbf{Registro actual de la labor}} \\ \hline
En papel & 42 & 67{,}7 & Sustenta \id{RF-04}, \id{RF-28} \\ \hline
Lo anota la administración & 13 & 21{,}0 & Sustenta \id{RF-35} \\ \hline
No se registra & 6 & 9{,}7 & Pérdida de información \\ \hline
\multicolumn{4}{|l|}{\textbf{Mayor dificultad declarada} (selección múltiple)} \\ \hline
Detectar plagas o enfermedades a tiempo & 25 & 40{,}3 & Sustenta \id{RF-07}, \id{RF-12} \\ \hline
El clima & 21 & 33{,}9 & Sustenta \id{RF-17} \\ \hline
Llevar la cuenta de lo realizado & 13 & 21{,}0 & Sustenta \id{RF-28}, \id{RF-29} \\ \hline
Saber qué producto aplicar & 13 & 21{,}0 & Sustenta \id{RF-09} \\ \hline
Comunicación con la administración & 6 & 9{,}7 & Sustenta \id{RF-30} \\ \hline
\multicolumn{4}{|l|}{\textbf{Capacidad tecnológica}} \\ \hline
Usa teléfono inteligente & 55 & 88{,}7 & \textbf{11{,}3\% no lo usa}: \id{RD-10} \\ \hline
Comodidad con aplicaciones (media 1--5) & 3{,}14 & --- & Sustenta \id{RNF-08}, \id{RD-06} \\ \hline
Señal de internet permanente & 16 & 25{,}8 & \textbf{74{,}2\% sin señal permanente}: \id{RD-02} \\ \hline
\multicolumn{4}{|l|}{\textbf{Aceptación}} \\ \hline
No le preocupa el registro por GPS & 46 & 74{,}2 & \textbf{25{,}8\% con reservas}: \id{RL-03} \\ \hline
Usaría la aplicación & 53 & 85{,}5 & Adopción esperada \\ \hline
\caption{Resultados del cuestionario ampliado ($n=62$, \id{EV-12})}
\end{longtable}

\textbf{Utilidad percibida de las funciones propuestas.} Cada persona respondiente valoró cinco
funciones en escala de 1 a 5. El orden resultante confirma la priorización MoSCoW de la §5.

\begin{longtable}{|L{7.6cm}|C{1.8cm}|C{1.8cm}|C{2.6cm}|}
\hline
\rowcolor{grisclaro}\textbf{Función propuesta} & \textbf{Media} & \textbf{Mediana} & \textbf{Requisito} \\
\hline
\endhead
Recibir alertas ante un problema del cultivo & \textbf{4{,}31} & 5 & \id{RF-12} \\ \hline
Foto $\rightarrow$ detectar plaga o enfermedad & \textbf{4{,}27} & 5 & \id{RF-07} \\ \hline
Foto $\rightarrow$ detectar deficiencia nutricional & 4{,}10 & 4 & \id{RF-08} \\ \hline
Registrar labores en el teléfono & 3{,}95 & 4 & \id{RF-04}, \id{RF-28} \\ \hline
Conteo automático de flores por GPS & 3{,}94 & 4 & \id{RF-14} \\ \hline
\caption{Utilidad percibida de las funciones propuestas ($n=62$, escala 1--5)}
\end{longtable}

\noindent
Las dos funciones mejor valoradas ---alertas y detección de plagas por imagen--- coinciden con la
dificultad más declarada por el personal (detectar plagas a tiempo, 40{,}3\%) y con la necesidad
expresada de forma independiente por la administración general (\id{EV-01}) y la asesoría técnica
(\id{EV-02}). Esta convergencia entre tres fuentes distintas ---cuestionario, entrevista operativa y
entrevista directiva--- sustenta la clasificación \emph{Must} de \id{RF-07} y \id{RF-12}.

\textbf{Datos faltantes.} La pregunta sobre comodidad con aplicaciones registró 37 respuestas
válidas de 62 ($40{,}3\%$ de no respuesta). Este valor se reporta como limitación y no se emplea
para inferencia; la media de 3{,}14 se presenta únicamente como indicio descriptivo.

% =====================================================================
%  3. REQUISITOS ESPECIFICOS
% =====================================================================
\newpage
\section{Requisitos específicos}

Esta sección especifica 42 requisitos funcionales, 19 requisitos no funcionales cuantificados sobre
las nueve características de calidad de ISO/IEC 25010:2023 \cite{iso25010}, 10 restricciones de
diseño y el mapeo de requisitos legales a la LOPDP \cite{lopdp}.

\subsection{Interfaces externas}

\subsubsection{Interfaces de usuario}

Las pantallas principales se especifican mediante prototipos de alta fidelidad identificados
\id{MU-01} a \id{MU-08}, incrustados en el Apéndice F y disponibles como archivos PNG en
\id{03\_Modelado/Mockups/}. La Tabla~\ref{tab:mockups} traza cada prototipo con los requisitos que
soporta.

\begin{longtable}{|C{1.6cm}|L{5.6cm}|L{7.6cm}|}
\hline
\rowcolor{grisclaro}\textbf{ID} & \textbf{Pantalla} & \textbf{Requisitos soportados} \\
\hline
\endhead
\id{MU-01} & Inicio de sesión & \id{RF-01} \\ \hline
\id{MU-02} & Panel principal & \id{RF-12}, \id{RF-19}, \id{RF-26}, \id{RF-29}, \id{RF-36} a \id{RF-39} \\ \hline
\id{MU-03} & Captura y análisis de imagen (IA) & \id{RF-07}, \id{RF-08}, \id{RF-09}, \id{RF-21} \\ \hline
\id{MU-04} & Registro de labor y polinización & \id{RF-04}, \id{RF-13}, \id{RF-14}, \id{RF-28} \\ \hline
\id{MU-05} & Detalle de lote & \id{RF-02}, \id{RF-05}, \id{RF-06}, \id{RF-11} \\ \hline
\id{MU-06} & Mapa de recorridos y polinización (GPS) & \id{RF-14}, \id{RF-15}, \id{RF-30} \\ \hline
\id{MU-07} & Reportes y estimación & \id{RF-18}, \id{RF-19}, \id{RF-32} \\ \hline
\id{MU-08} & Alertas & \id{RF-12}, \id{RF-22} \\ \hline
\caption{Trazabilidad entre prototipos de interfaz y requisitos funcionales}
\label{tab:mockups}
\end{longtable}

\subsubsection{Interfaces de hardware}

Cámara del dispositivo móvil (captura de imagen para los tres componentes de IA); receptor GPS
(conteo georreferenciado y registro de recorridos); dispositivo Android 9.0 o superior; servidor o
servicio en la nube para almacenamiento e inferencia.

\subsubsection{Interfaces de software}

API REST sobre HTTPS con carga útil JSON entre cliente móvil, panel web y backend; servicio de
inferencia de IA; servicio de autenticación con control de acceso basado en roles; base de datos
relacional; servicio de datos climáticos.

\subsubsection{Interfaces de comunicación}

Se incorpora en esta entrega la \textbf{interfaz de intercambio con la planta extractora}. La
evidencia \id{EV-04} documenta que la extractora opera un sistema interno propio que emite un ticket
de pesaje con peso bruto, tara, neto y calificación de calidad, y que lo remite al productor por
correo electrónico. El SIMPA no se integra directamente con ese sistema en la versión v1; el
intercambio se realiza mediante exportación de archivos en formato estructurado, conforme a
\id{RF-32} y a la restricción de minimización de datos \id{RD-09}.

\subsection{Requisitos funcionales}

Cada requisito se documenta con la plantilla de ocho atributos e incluye la referencia explícita al
identificador de evidencia que lo respalda, conforme a la §3.3 de la guía entonces aplicable a la Entrega 3 (2A).

\subsubsection{Bloque I --- Gestión de la estructura productiva}

\RF{RF-01}{Autenticación y control de acceso por rol}
{El sistema debe permitir el inicio de sesión y restringir las funciones disponibles según el rol de la persona usuaria (administrador general, administrador/asesor técnico, asistente de administración, jefe de polinización, trabajador agrícola).}
{Administrador general · \id{EV-01}, \id{EV-07}, \id{EV-09}}
{\textbf{Entrada:} credenciales de acceso. \textbf{Salida:} sesión activa con menú correspondiente al rol, o mensaje de error.}
{\textbf{Pre:} la persona usuaria está registrada y activa. \textbf{Post:} accede únicamente a las funciones autorizadas para su rol.}
{Must}
{Una cuenta con rol \emph{trabajador agrícola} no puede acceder a ninguna función de administración; el acceso se concede exclusivamente con credenciales válidas.}

\RF{RF-02}{Gestión de plantaciones y lotes}
{El sistema debe permitir registrar, editar y consultar la plantación y sus lotes, con variedad sembrada, área en hectáreas, número de palmas y fecha de siembra.}
{Administrador general · \id{EV-01}, \id{EV-02}, \id{EV-11}}
{\textbf{Entrada:} datos identificativos y productivos del lote. \textbf{Salida:} lote almacenado y disponible en el listado.}
{\textbf{Pre:} persona usuaria con permiso de administración. \textbf{Post:} el lote queda disponible para asociar labores, diagnósticos y cosechas.}
{Must}
{Tras registrar un lote, este aparece en el listado y puede seleccionarse como destino de una labor, un diagnóstico y un registro de cosecha.}

\RF{RF-03}{Gestión de personal y equipos de trabajo}
{El sistema debe permitir registrar al personal y organizarlo por equipos funcionales: administración, polinización, control fitosanitario, riego y labores agrícolas.}
{Administrador general · \id{EV-01}, \id{EV-11}}
{\textbf{Entrada:} datos laborales de la persona trabajadora y equipo asignado. \textbf{Salida:} persona registrada y vinculada a un equipo.}
{\textbf{Pre:} persona usuaria con permiso de administración. \textbf{Post:} queda disponible para asignación de labores y cálculo de avance.}
{Must}
{Cada persona registrada queda asociada al menos a un equipo y aparece como asignable en el registro de labores.}

\RF{RF-10}{Gestión de variedades de palma}
{El sistema debe registrar la variedad sembrada por lote y emplearla para parametrizar los umbrales de diagnóstico y de madurez, así como las advertencias sobre polinizadores naturales.}
{Administrador / Asesor técnico · \id{EV-01}, \id{EV-02}, \id{EV-04}}
{\textbf{Entrada:} variedad asignada al lote y sus parámetros. \textbf{Salida:} umbrales aplicables a los módulos de diagnóstico y clasificación.}
{\textbf{Pre:} el lote existe. \textbf{Post:} los módulos de IA operan con los umbrales de la variedad correspondiente.}
{Must}
{Para un lote de variedad \emph{guineensis}, el umbral de madurez aplicado es de 3 frutos desprendidos; para un lote de variedad híbrida, de 5.}

\subsubsection{Bloque II --- Planificación y seguimiento de labores}

\RF{RF-26}{Planificación semanal de labores con presupuesto}
{El sistema debe permitir registrar el presupuesto semanal de labores por lote, expresado en la unidad de avance correspondiente, y contrastarlo con el avance efectivamente cumplido al cierre de la semana.}
{Administrador general · \id{EV-11}}
{\textbf{Entrada:} semana, lote, tipo de labor y meta en unidades. \textbf{Salida:} plan semanal y reporte comparativo presupuestado \emph{versus} cumplido.}
{\textbf{Pre:} los lotes y los tipos de labor están registrados. \textbf{Post:} el plan queda disponible para el seguimiento diario y el cierre semanal.}
{Must}
{Al cierre de una semana, el sistema muestra para cada labor planificada la meta, el avance cumplido y el porcentaje de cumplimiento.}

\RF{RF-27}{Arrastre de labor incompleta a la semana siguiente}
{Cuando una labor planificada no alcanza su meta al cierre de la semana, el sistema debe generar automáticamente el remanente como pendiente en el plan de la semana siguiente, conservando la referencia a la labor de origen.}
{Administrador general · \id{EV-11}}
{\textbf{Entrada:} cierre de la semana con avance inferior a la meta. \textbf{Salida:} labor pendiente creada en la semana siguiente con la cantidad remanente.}
{\textbf{Pre:} existe un plan semanal cerrado con cumplimiento parcial. \textbf{Post:} el remanente aparece en el plan siguiente marcado como continuación.}
{Should}
{Dada una meta de 1.123 coronas con 800 cumplidas, el sistema genera en la semana siguiente una labor pendiente de 323 coronas vinculada a la original.}

\RF{RF-34}{Registro de solicitud de insumos asociada al plan}
{El sistema debe permitir registrar los insumos requeridos para ejecutar el plan semanal y su estado de provisión.}
{Administrador general · \id{EV-11}}
{\textbf{Entrada:} insumo, cantidad y semana asociada. \textbf{Salida:} listado de requerimientos por semana con su estado.}
{\textbf{Pre:} existe un plan semanal. \textbf{Post:} el requerimiento queda asociado a la semana y consultable por la administración.}
{Could}
{Un insumo registrado como requerimiento de la semana aparece en el listado de la semana correspondiente con su estado de provisión.}

\RF{RF-36}{Catálogo codificado de labores con tarifa por unidad}
{El sistema debe mantener un catálogo de tipos de labor con su código identificador, su unidad de medida asociada (jornal, hora, avance o tonelada) y su tarifa vigente, de modo que la tarifa aplicable a un registro de avance se obtenga del catálogo y no se ingrese de forma manual en cada registro.}
{Administrador general · \id{EV-13}}
{\textbf{Entrada:} código de labor, unidad de medida, tarifa y fecha de vigencia. \textbf{Salida:} catálogo consultable y tarifa aplicable a cada registro de avance.}
{\textbf{Pre:} persona usuaria con permiso de administración. \textbf{Post:} los registros de avance posteriores emplean la tarifa vigente a su fecha.}
{Must}
{Registrada una labor de tipo corona con tarifa por unidad de avance, el cálculo del valor de un registro de 400 unidades corresponde al producto entre la cantidad y la tarifa vigente en la fecha del registro.}

\RF{RF-37}{Liquidación semanal de presupuesto contra ejecución}
{El sistema debe generar, al cierre de cada semana, una liquidación que contraste por grupo de labor la cantidad presupuestada, la cantidad ejecutada y el valor resultante, y que totalice el gasto presupuestado frente al gasto real.}
{Administrador general · \id{EV-13}, \id{EV-11}}
{\textbf{Entrada:} plan semanal y registros de avance del período. \textbf{Salida:} liquidación por grupo con totales y diferencia entre presupuestado y real.}
{\textbf{Pre:} existe un plan semanal y registros de avance asociados. \textbf{Post:} la liquidación queda disponible para aprobación.}
{Must}
{Para una semana con plan y avances registrados, el sistema presenta por cada grupo de labor la cantidad, la tarifa, el valor y el total, y calcula la diferencia entre el gasto presupuestado y el real.}

\RF{RF-38}{Registro de labor no presupuestada}
{El sistema debe permitir registrar labores ejecutadas que no figuraban en el plan de la semana, identificándolas de forma diferenciada para que no se confundan con el cumplimiento de las metas planificadas.}
{Administrador general · \id{EV-13}}
{\textbf{Entrada:} labor, cantidad, lote y motivo. \textbf{Salida:} registro marcado como no presupuestado, excluido del cálculo de cumplimiento del plan.}
{\textbf{Pre:} existe un plan semanal publicado. \textbf{Post:} la labor computa en el gasto real pero no en el cumplimiento de las metas.}
{Should}
{Una labor registrada como no presupuestada aparece en el grupo correspondiente de la liquidación y no altera el porcentaje de cumplimiento del plan.}

\RF{RF-39}{Aprobación de la liquidación semanal}
{El sistema debe requerir la aprobación explícita de la persona administradora para cerrar la liquidación de una semana, dejando constancia de quién aprobó y en qué fecha, tras lo cual los registros del período no admiten modificación sin trazabilidad.}
{Administrador general · \id{EV-13}}
{\textbf{Entrada:} liquidación semanal completa. \textbf{Salida:} liquidación cerrada con constancia de aprobación.}
{\textbf{Pre:} la liquidación de la semana está generada. \textbf{Post:} el período queda cerrado; toda modificación posterior genera entrada en bitácora.}
{Should}
{Aprobada una liquidación, el sistema registra la identidad de quien aprueba y la fecha, e impide modificar los registros del período sin dejar constancia.}


\RF{RF-35}{Registro delegado para personal sin dispositivo propio}
{El sistema debe permitir que una persona con rol de supervisión registre el avance de labor en nombre de una persona trabajadora que no dispone de teléfono inteligente, dejando constancia de quién efectuó el registro y en nombre de quién.}
{Jefe de campo / Asistente de administración · \id{EV-12}, \id{EV-05}}
{\textbf{Entrada:} identificación de la persona trabajadora, labor, cantidad y unidad. \textbf{Salida:} avance registrado a nombre de la persona trabajadora, con marca de registro delegado.}
{\textbf{Pre:} la persona registradora tiene rol de supervisión y la persona trabajadora está registrada. \textbf{Post:} el avance computa para la persona trabajadora y la bitácora conserva la identidad de quien lo registró.}
{Must}
{Un avance registrado de forma delegada aparece atribuido a la persona trabajadora en su consolidado, y la bitácora identifica a la persona supervisora que lo ingresó.}

\subsubsection{Bloque III --- Registro operativo de campo}

\RF{RF-04}{Registro de labores agrícolas}
{El sistema debe permitir registrar las labores ejecutadas por lote: chapia, corona, control químico de malezas, eliminación de flores andróginas, sanidad vegetal, poda, riego, polinización y cosecha.}
{Trabajador agrícola / Jefe de campo · \id{EV-01}, \id{EV-02}, \id{EV-05}, \id{EV-06}, \id{EV-11}}
{\textbf{Entrada:} tipo de labor, lote, fecha y persona responsable. \textbf{Salida:} labor registrada y asociada al lote, la fecha y la persona.}
{\textbf{Pre:} el lote y la persona trabajadora existen. \textbf{Post:} la labor queda en el historial del lote y computa para el avance.}
{Must}
{Una labor registrada queda vinculada a su lote, fecha y responsable, y es recuperable mediante consulta del historial del lote.}

\RF{RF-28}{Registro de avance por unidad de labor}
{El sistema debe registrar el avance diario de cada persona trabajadora en la unidad propia de cada labor: coronas, bancos, plantas, matas, jornales, racimos o toneladas, según corresponda al tipo de labor.}
{Trabajador agrícola · \id{EV-05}, \id{EV-06}, \id{EV-11}}
{\textbf{Entrada:} labor, cantidad ejecutada y unidad. \textbf{Salida:} avance diario acumulado por persona, lote y semana.}
{\textbf{Pre:} existe una labor asignada. \textbf{Post:} el avance queda registrado y disponible para el cálculo de cumplimiento y de remuneración por avance.}
{Must}
{Para una jornada de chapia, el sistema registra el número de matas ejecutadas y lo acumula al total semanal de la persona trabajadora.}

\RF{RF-29}{Consulta de avance acumulado y estimación de remuneración}
{El sistema debe permitir a la persona trabajadora consultar en su propio dispositivo el avance acumulado de la quincena y la estimación de remuneración correspondiente al trabajo por avance.}
{Trabajador agrícola · \id{EV-06}, \id{EV-07}}
{\textbf{Entrada:} identificación de la persona trabajadora y período. \textbf{Salida:} avance acumulado por labor y estimación de remuneración del período.}
{\textbf{Pre:} existen registros de avance en el período. \textbf{Post:} se muestra el consolidado sin posibilidad de edición por parte de la persona trabajadora.}
{Should}
{Una persona trabajadora consulta su avance quincenal y obtiene el desglose por labor y el valor estimado, coincidente con la suma de sus registros diarios.}

\RF{RF-05}{Registro de monitoreo fitosanitario}
{El sistema debe permitir registrar el monitoreo del estado fitosanitario por lote y marcar las incidencias detectadas.}
{Técnico fitosanitario · \id{EV-01}, \id{EV-02}, \id{EV-09}}
{\textbf{Entrada:} lote, observaciones y estado. \textbf{Salida:} registro de monitoreo y, si procede, alerta asociada.}
{\textbf{Pre:} persona usuaria autenticada. \textbf{Post:} el monitoreo queda disponible para consulta y reportes.}
{Must}
{El sistema almacena cada monitoreo con fecha y lote, y permite filtrar los lotes que presentan incidencia abierta.}

\RF{RF-30}{Reporte de incidencia desde campo con evidencia fotográfica y georreferencia}
{El sistema debe permitir reportar una afectación desde el campo adjuntando fotografía y coordenadas del punto exacto, de modo que la verificación posterior no dependa del recuerdo de la persona que reportó.}
{Asistente de administración · \id{EV-07}}
{\textbf{Entrada:} fotografía, lote, coordenadas y descripción. \textbf{Salida:} incidencia registrada y georreferenciada, visible en el mapa del lote.}
{\textbf{Pre:} persona usuaria autenticada con GPS disponible. \textbf{Post:} la incidencia queda pendiente de verificación con su ubicación exacta.}
{Must}
{Una incidencia reportada en campo es recuperable posteriormente con su fotografía y sus coordenadas, y su ubicación se muestra sobre el mapa del lote.}

\RF{RF-31}{Lista de verificación de pendientes por visita}
{El sistema debe mantener, para cada persona con rol de supervisión, una lista de actividades pendientes de verificación en su próxima visita a campo, alimentada por las incidencias reportadas.}
{Asistente de administración · \id{EV-07}}
{\textbf{Entrada:} incidencias reportadas y no verificadas. \textbf{Salida:} lista de pendientes ordenada por antigüedad y criticidad.}
{\textbf{Pre:} existen incidencias sin verificar. \textbf{Post:} la lista se actualiza al marcar cada pendiente como verificado.}
{Should}
{Al reportarse una incidencia, esta aparece en la lista de pendientes de la persona supervisora y desaparece al marcarse como verificada.}

\RF{RF-06}{Seguimiento de etapas del cultivo}
{El sistema debe registrar y dar seguimiento a la etapa del cultivo por lote: siembra, mantenimiento, polinización y cosecha.}
{Administrador general · \id{EV-01}}
{\textbf{Entrada:} etapa y fecha de transición por lote. \textbf{Salida:} línea de tiempo de etapas del lote.}
{\textbf{Pre:} el lote existe. \textbf{Post:} la etapa actual queda actualizada.}
{Should}
{Cada lote muestra su etapa actual y el historial completo de transiciones con sus fechas.}

\RF{RF-40}{Exportación de los datos personales de la persona titular}
{El sistema debe permitir a cada persona trabajadora obtener una copia de los datos personales que el sistema conserva sobre ella --- datos identificativos, de contacto, de vinculación laboral y de avance registrado --- en un formato estructurado y de uso común, sin intermediación de la administración.}
{Trabajador agrícola · \id{RL-04} (Art. 12 LOPDP)}
{\textbf{Entrada:} solicitud de la persona titular autenticada. \textbf{Salida:} archivo en formato CSV o PDF con la totalidad de sus datos personales conservados.}
{\textbf{Pre:} sesión iniciada por la propia persona titular. \textbf{Post:} la solicitud y su fecha quedan registradas en la bitácora de accesos.}
{Must}
{Solicitada la exportación por una persona titular, el archivo generado contiene la totalidad de los campos personales que el sistema conserva sobre ella y ninguno correspondiente a otra persona. La solicitud figura en la bitácora con su marca temporal.}

\RF{RF-41}{Rectificación de datos con conservación en bitácora}
{El sistema debe permitir corregir los datos personales y los registros de labor de una persona trabajadora conservando el valor anterior, la fecha, la identificación de quien corrige y el motivo declarado, sin que la corrección elimine la evidencia del valor sustituido.}
{Administrador general / Trabajador agrícola · \id{RL-05} (Art. 13 LOPDP)}
{\textbf{Entrada:} registro a corregir, valor nuevo y motivo. \textbf{Salida:} registro actualizado y entrada de bitácora con el valor anterior.}
{\textbf{Pre:} existe el registro y la persona que corrige tiene rol autorizado. \textbf{Post:} el valor anterior permanece consultable en la bitácora de forma indefinida.}
{Must}
{Corregido un dato, la consulta del registro devuelve el valor nuevo y la consulta de la bitácora devuelve el valor anterior con fecha, autor y motivo. Ninguna operación de la interfaz permite borrar una entrada de bitácora.}

\RF{RF-42}{Supresión de datos personales con disociación del histórico productivo}
{El sistema debe suprimir los datos identificativos y de contacto de una persona trabajadora al término de su relación laboral, sustituyendo su identificador por uno disociado que no permita la reidentificación, y conservando los registros de avance y producción asociados a ese identificador para el histórico de los lotes.}
{Administrador general · \id{RL-06} (Art. 14 LOPDP)}
{\textbf{Entrada:} relación laboral marcada como terminada y confirmación de la administración. \textbf{Salida:} datos personales suprimidos e histórico productivo conservado de forma disociada.}
{\textbf{Pre:} la relación laboral consta como terminada. \textbf{Post:} los datos identificativos no son recuperables por ninguna función del sistema; los registros productivos permanecen agregables por lote y por período.}
{Must}
{Ejecutada la supresión, ninguna consulta ni exportación del sistema devuelve el nombre, la cédula ni el contacto de la persona; los totales de avance del lote y del período permanecen idénticos a los previos a la supresión.}

\RF{RF-11}{Registro de análisis de suelo y foliar}
{El sistema debe registrar los resultados de los análisis de suelo y foliar por lote, incluyendo nitrógeno, fósforo, potasio y elementos menores, para guiar la fertilización.}
{Administrador / Asesor técnico · \id{EV-02}, \id{EV-04}}
{\textbf{Entrada:} valores del análisis por lote y fecha. \textbf{Salida:} registro histórico y comparativo por lote.}
{\textbf{Pre:} el lote existe. \textbf{Post:} el análisis queda asociado al lote y consultable por fecha.}
{Should}
{Tras registrar un análisis, sus valores quedan asociados al lote y pueden compararse con los análisis previos del mismo lote.}

\RF{RF-17}{Registro y consulta de datos climáticos}
{El sistema debe registrar y permitir consultar los datos climáticos del cultivo: milímetros de lluvia medidos con pluviómetro y datos de estaciones meteorológicas.}
{Administrador general / Asesor técnico · \id{EV-01}, \id{EV-02}, \id{EV-03}}
{\textbf{Entrada:} lectura de precipitación y fecha. \textbf{Salida:} histórico climático y consulta de lluvia acumulada por período.}
{\textbf{Pre:} persona usuaria autenticada. \textbf{Post:} la lectura queda incorporada al registro climático del cultivo.}
{Should}
{El sistema permite consultar la precipitación acumulada en un rango de fechas definido por la persona usuaria.}

\subsubsection{Bloque IV --- Polinización asistida}

\RF{RF-13}{Registro del proceso de polinización}
{El sistema debe registrar la polinización por lote, incluyendo el estado de la antesis (coloración amarilla o negra) y la verificación visual de la aplicación del polinizante.}
{Jefe de polinización / Polinizador · \id{EV-03}}
{\textbf{Entrada:} lote, estado de antesis y verificación visual. \textbf{Salida:} registro de polinización asociado al lote, fecha y persona.}
{\textbf{Pre:} el lote se encuentra en etapa de polinización. \textbf{Post:} el registro queda disponible para el control de calidad de la labor.}
{Must}
{El sistema permite registrar la polinización indicando el estado de antesis y consultar los lotes polinizados en una fecha dada.}

\RF{RF-14}{Conteo georreferenciado de flores polinizadas}
{El sistema debe calcular automáticamente, a partir de las marcaciones georreferenciadas registradas durante la jornada, el número de flores polinizadas por cada persona trabajadora, sin requerir el ingreso manual planta por planta.}
{Polinizador · \id{EV-03}}
{\textbf{Entrada:} marcaciones georreferenciadas de la jornada. \textbf{Salida:} total diario de flores polinizadas por persona y lote.}
{\textbf{Pre:} sesión iniciada con rastreo GPS activo. \textbf{Post:} el conteo queda registrado por línea de siembra recorrida, no por planta individual. \textbf{Nota de alcance:} \id{RNF-05} admite una precisión GPS de hasta 10 m. El ERS no declara la distancia entre palmas de la plantación, por lo que no es posible afirmar que esa precisión permita discriminar una palma de la contigua; el marco de siembra declarado por la organización es de 9,50~m. En consecuencia, esta versión acota el requisito a la agregación por línea de siembra y excluye de su alcance la asignación individual planta por planta. La verificación de si el marco de siembra permite el conteo individual se registra como \id{L-12}.}
{Must}
{Al cierre de la jornada, el sistema calcula el total de flores polinizadas por persona a partir de las marcaciones GPS, sin registro manual individual.}

\RF{RF-15}{Registro y visualización de recorridos del personal}
{El sistema debe registrar y mostrar los recorridos georreferenciados del personal en campo, con el fin de verificar la cobertura efectiva de las labores asignadas.}
{Administrador general · \id{EV-01}, \id{EV-02}, \id{EV-03}}
{\textbf{Entrada:} secuencia de ubicaciones durante la jornada. \textbf{Salida:} recorrido representado sobre el mapa del lote, por persona y fecha.}
{\textbf{Pre:} la persona trabajadora ha otorgado consentimiento para el tratamiento de datos de geolocalización (véase \id{RL-03}). \textbf{Post:} el recorrido queda disponible para supervisión.}
{Should}
{La administración puede visualizar el recorrido de una persona trabajadora en una fecha determinada, siempre que exista consentimiento vigente.}

\RF{RF-16}{Evaluación de desempeño del personal}
{El sistema debe calcular indicadores de desempeño por persona trabajadora a partir del avance registrado, la cobertura de los recorridos y la calidad resultante de la fruta cosechada en los lotes bajo su responsabilidad.}
{Administrador general · \id{EV-01}, \id{EV-02}, \id{EV-03}}
{\textbf{Entrada:} registros de avance, recorridos y calificaciones de calidad del período. \textbf{Salida:} indicador de desempeño por persona y período, expresado en escala de 0 a 100, calculado como el promedio ponderado de tres componentes: cumplimiento de la meta de avance (peso 0,5), cobertura del recorrido planificado (peso 0,3) y calificación de calidad de los lotes bajo responsabilidad (peso 0,2), cada componente normalizado a [0, 100]. Los pesos son configurables por la administración.}
{\textbf{Pre:} existen al menos cinco jornadas con registro de avance en el período. \textbf{Post:} el indicador queda disponible para consulta de la administración, junto con el desglose de sus tres componentes.}
{Should}
{Dado un período con registros de avance, recorrido y calidad conocidos, el indicador calculado por el sistema coincide con el promedio ponderado de los tres componentes con una tolerancia de $\pm 1$ punto, y el desglose presentado reproduce los valores de entrada.}

\subsubsection{Bloque V --- Diagnóstico asistido por inteligencia artificial}

\noindent\textbf{Viabilidad técnica de los componentes de IA.} Los tres requisitos de análisis de
imagen de este bloque no son especulativos: la literatura reciente documenta implementaciones
equivalentes en el mismo cultivo. La clasificación de madurez de racimos de palma aceitera
\emph{sobre dispositivos móviles} mediante aprendizaje profundo ha sido reportada por Suharjito
\emph{et al.} \cite{suharjito2021}, y la revisión de Goh \emph{et al.} \cite{goh2025} sistematiza
los métodos disponibles y sus tasas de acierto. Existen además conjuntos de datos públicos de
racimos capturados en exteriores \cite{omar2024} y aproximaciones basadas en segmentación y
características texturales \cite{rosbi2024} que evitan la dependencia del color, consistente con la
restricción \id{RD-07}. Rizzo \emph{et al.} \cite{rizzo2023} ofrecen el panorama general de la
clasificación de madurez de frutos, y Kipli \emph{et al.} \cite{kipli2023} documentan aplicaciones
de detección y conteo de palmas. Zaki \emph{et al.} \cite{zaki2025} sitúan estas tecnologías dentro
de la transformación digital del sector palmicultor. Estas fuentes sustentan los valores objetivo
fijados en \id{RNF-01} y \id{RNF-02}.

\vspace{0.3cm}
\RF{RF-07}{Detección de plagas y enfermedades por análisis de imagen}
{El sistema debe analizar una fotografía de la palma e identificar la presencia de plagas o enfermedades, indicando el tipo probable de afectación.}
{Técnico fitosanitario / Trabajador agrícola · \id{EV-01}, \id{EV-02}, \id{EV-08}}
{\textbf{Entrada:} fotografía de la palma o de la hoja. \textbf{Salida:} clasificación de la afectación con nivel de confianza y explicación asociada (véase \id{RNF-16}).}
{\textbf{Pre:} la imagen cumple los requisitos mínimos de calidad, que esta ficha define como: resolución no inferior a 1024$\times$768 px, formato JPEG o PNG, tamaño no superior a 5 MB y encuadre en el que el órgano vegetal afectado ocupe al menos un tercio del área. \textbf{Post:} el diagnóstico queda asociado al lote y a la fecha; si la imagen no cumple alguna condición, el sistema la rechaza indicando cuál.}
{Must}
{Ante un conjunto de imágenes de prueba con afectación previamente conocida, el sistema devuelve la clasificación correspondiente y la registra con su lote y fecha.}

\RF{RF-08}{Diagnóstico de deficiencias nutricionales por análisis de imagen}
{El sistema debe analizar la fotografía de una hoja y estimar la deficiencia nutricional presente, sugiriendo la corrección correspondiente.}
{Administrador / Asesor técnico · \id{EV-02}, \id{EV-04}}
{\textbf{Entrada:} fotografía de la hoja. \textbf{Salida:} deficiencia estimada, recomendación de fertilización y explicación asociada.}
{\textbf{Pre:} imagen válida y lote con variedad asignada. \textbf{Post:} el diagnóstico nutricional queda registrado en el historial del lote.}
{Must}
{Ante una hoja con deficiencia previamente conocida, el sistema indica la deficiencia, propone una corrección y almacena el registro.}

\RF{RF-09}{Ficha técnica de plaga con recomendación de control}
{El sistema debe presentar, para cada plaga detectada, su especie o familia, su ciclo de evolución y la acción de control recomendada según la etapa del insecto.}
{Técnico fitosanitario · \id{EV-01}, \id{EV-02}}
{\textbf{Entrada:} plaga identificada por \id{RF-07}. \textbf{Salida:} ficha con especie, ciclo y acción recomendada diferenciada por etapa.}
{\textbf{Pre:} existe un diagnóstico de plaga. \textbf{Post:} la recomendación queda disponible y registrada junto al diagnóstico.}
{Should}
{Para una plaga detectada en estado adulto, el sistema recomienda trampeo; para la misma plaga en estado larval, recomienda aplicación de insecticida.}

\RF{RF-21}{Clasificación de madurez del racimo por desprendimiento}
{El sistema debe estimar el estado de madurez de un racimo a partir de una fotografía de su parte basal, empleando el conteo de frutos desprendidos como criterio y aplicando el umbral correspondiente a la variedad del lote. El sistema no debe emplear el color del fruto como criterio de clasificación.}
{Técnico de la extractora / Cosechador · \id{EV-04}, \id{EV-08}}
{\textbf{Entrada:} fotografía de la parte basal del racimo y variedad del lote. \textbf{Salida:} estado estimado (verde, maduro, sobremaduro) con nivel de confianza y explicación.}
{\textbf{Pre:} el lote tiene variedad asignada (\id{RF-10}). \textbf{Post:} la estimación queda asociada al racimo, el lote y la fecha.}
{Must}
{Para un lote de variedad híbrida, un racimo con menos de 5 frutos desprendidos se clasifica como verde; con 5 o más, como maduro. Para variedad \emph{guineensis} el umbral aplicado es de 3.}

\RF{RF-12}{Generación de alertas tempranas}
{El sistema debe generar alertas ante cualquiera de los cuatro disparadores siguientes, dirigidas al rol responsable de su atención: (a) diagnóstico positivo de plaga o enfermedad (\id{RF-07}); (b) diagnóstico positivo de deficiencia nutricional (\id{RF-08}); (c) porcentaje estimado de fruta verde que alcance o supere el umbral configurado (\id{RF-22}); (d) ausencia de registro de avance en un lote con labor planificada para la jornada (\id{RF-26}, \id{RF-28}). No se admiten disparadores distintos de los enumerados.}
{Técnico / Administrador general · \id{EV-01}, \id{EV-02}}
{\textbf{Entrada:} diagnóstico o incidencia registrada. \textbf{Salida:} alerta con lote, tipo, criticidad y fecha, visible para el rol correspondiente.}
{\textbf{Pre:} se cumple al menos uno de los cuatro disparadores enumerados en la descripción. \textbf{Post:} la alerta queda registrada con estado abierto hasta su atención.}
{Must}
{Provocado cada uno de los cuatro disparadores por separado, el sistema genera en los cuatro casos una alerta asociada al lote, visible para el rol responsable, con tipo, criticidad, estado y fecha. Provocada una condición no enumerada, no se genera alerta.}

\subsubsection{Bloque VI --- Calidad de la fruta y relación con la extractora}

\RF{RF-22}{Estimación y alerta preventiva de porcentaje de fruta verde}
{El sistema debe estimar, antes del despacho de un lote cosechado, el porcentaje de racimos clasificados como verdes, y emitir una alerta cuando dicho porcentaje alcance el 80\% del umbral a partir del cual la planta extractora aplica penalización. Con el umbral de penalización configurado en el 3\%, la alerta preventiva se dispara a partir del 2,4\%.}
{Administrador general / Técnico de la extractora · \id{EV-04}, \id{EV-08}}
{\textbf{Entrada:} clasificaciones de madurez del lote cosechado (\id{RF-21}). \textbf{Salida:} porcentaje estimado de fruta verde y alerta si supera el umbral configurado.}
{\textbf{Pre:} existen racimos clasificados en el lote. \textbf{Post:} la alerta queda registrada antes del despacho.}
{Must}
{Configurado un umbral de penalización del 3\%, un lote con estimación de fruta verde del 2,4\% o superior genera la alerta preventiva antes del despacho, y un lote con estimación del 2,3\% no la genera.}

\RF{RF-23}{Registro de la calificación de calidad emitida por la extractora}
{El sistema debe permitir registrar los datos del ticket de pesaje emitido por la planta extractora: peso bruto, tara, peso neto y calificación de calidad (tamaño del fruto, fruta verde, sobremadura, malformada y pedúnculo largo).}
{Administrador general · \id{EV-04}}
{\textbf{Entrada:} datos del ticket de pesaje. \textbf{Salida:} registro de entrega asociado al lote y a la fecha, con su calificación.}
{\textbf{Pre:} existe un despacho registrado. \textbf{Post:} la calificación queda asociada al lote de origen y disponible para análisis histórico.}
{Should}
{Registrado un ticket, sus valores quedan asociados al lote de origen y son consultables en el histórico de entregas.}

\RF{RF-24}{Trazabilidad de la calificación hasta la cuadrilla responsable}
{El sistema debe permitir rastrear una calificación desfavorable de la extractora hasta el lote de origen, la fecha de corte y la cuadrilla de cosecha responsable, con el fin de distinguir si el problema proviene de la cosecha o del manejo del cultivo.}
{Técnico de la extractora / Administrador general · \id{EV-04}, \id{EV-08}}
{\textbf{Entrada:} identificación del despacho calificado. \textbf{Salida:} lote, fecha de corte, cuadrilla y labores asociadas al período.}
{\textbf{Pre:} el despacho tiene calificación registrada y el lote tiene labores registradas. \textbf{Post:} se muestra la cadena completa de trazabilidad.}
{Should}
{Dado un despacho con penalización por fruta verde, el sistema identifica el lote, la fecha de corte y la cuadrilla que ejecutó la cosecha.}

\RF{RF-25}{Control de la ventana entre corte y entrega}
{El sistema debe registrar la fecha y hora de corte de cada lote cosechado y calcular el tiempo transcurrido hasta la entrega en la planta extractora, emitiendo una advertencia cuando dicho intervalo supere el valor configurado para la variedad del lote. El valor es un atributo de la variedad, gestionado mediante \id{RF-10}, y no una constante del sistema: el ERS no fija su magnitud porque depende de la variedad sembrada y del acuerdo con la planta extractora. La administración lo configura por variedad antes de la puesta en operación.}
{Técnico de la extractora · \id{EV-04}}
{\textbf{Entrada:} marca temporal de corte y de entrega. \textbf{Salida:} intervalo calculado y advertencia si excede el umbral configurado.}
{\textbf{Pre:} existe un registro de cosecha con marca temporal y la variedad del lote tiene configurado su valor de ventana (\id{RF-10}). \textbf{Post:} el intervalo queda registrado junto al despacho.}
{Should}
{Configurado un umbral de 24 horas para la variedad del lote, un despacho entregado 30 horas después del corte genera la advertencia; uno entregado a las 23 horas no la genera. Si la variedad no tiene valor configurado, el sistema lo indica y no emite advertencia.}

\RF{RF-33}{Registro de la guía de remisión del despacho}
{El sistema debe permitir registrar el número y la fecha de la guía de remisión asociada a cada despacho de fruta, y advertir cuando un despacho se registre sin ella.}
{Técnico de la extractora · \id{EV-04}}
{\textbf{Entrada:} número y fecha de la guía de remisión. \textbf{Salida:} despacho con guía asociada, o advertencia de ausencia.}
{\textbf{Pre:} existe un despacho registrado. \textbf{Post:} la guía queda vinculada al despacho.}
{Could}
{Un despacho registrado sin número de guía de remisión genera una advertencia visible para la administración antes de su cierre.}

\subsubsection{Bloque VII --- Reportes, estimación e histórico}

\RF{RF-18}{Estimación de producción y rendimiento}
{El sistema debe estimar la producción futura a partir del censo de flores, el número de racimos en desarrollo y el rendimiento propio de la variedad sembrada.}
{Administrador general / Asesor técnico · \id{EV-01}, \id{EV-02}}
{\textbf{Entrada:} censo de flores, conteo de racimos y variedad. \textbf{Salida:} estimación de producción por lote y período.}
{\textbf{Pre:} existe al menos un censo registrado en los seis meses previos. \textbf{Post:} la estimación queda disponible para la planificación, acompañada de su margen de error declarado.}
{Must}
{Contrastada la estimación contra la producción real registrada del mismo lote y período, la desviación absoluta relativa no supera el 20\%. La estimación se presenta siempre acompañada de ese margen. Una desviación superior al 20\% se considera resultado incorrecto y debe registrarse como incidencia de calibración.}

\RF{RF-32}{Compartición de la estimación de cosecha con la planta extractora}
{El sistema debe permitir exportar la estimación de cosecha del período hacia la planta extractora, limitando la información compartida al conteo de racimos y al peso promedio estimado, con exclusión de cualquier dato personal o económico.}
{Técnico de la extractora · \id{EV-04}}
{\textbf{Entrada:} período y lotes seleccionados. \textbf{Salida:} archivo estructurado con conteo de racimos y peso promedio estimado.}
{\textbf{Pre:} existe una estimación calculada. \textbf{Post:} el archivo queda disponible para su envío, sin contener datos personales.}
{Should}
{El archivo exportado contiene exclusivamente conteo de racimos y peso promedio por período, y no incluye ningún campo de identificación personal ni información económica.}

\RF{RF-19}{Generación de reportes}
{El sistema debe generar reportes diarios, consolidados semanales y estadísticas anuales sobre el estado del cultivo, las labores ejecutadas, el personal y la producción.}
{Administrador general · \id{EV-01}, \id{EV-02}, \id{EV-11}}
{\textbf{Entrada:} tipo de reporte y período. \textbf{Salida:} reporte consolidado, exportable a PDF y a hoja de cálculo.}
{\textbf{Pre:} existen registros en el período seleccionado. \textbf{Post:} el reporte queda disponible para consulta o exportación.}
{Must}
{El sistema genera el consolidado semanal de labores y producción del período seleccionado y permite exportarlo.}

\RF{RF-20}{Historial de actividades, diagnósticos y eventos}
{El sistema debe mantener un historial consultable y filtrable de todas las actividades, diagnósticos, alertas y entregas asociadas a cada lote.}
{Administrador general / Técnico · \id{EV-01}}
{\textbf{Entrada:} filtros por lote, fecha y tipo de evento. \textbf{Salida:} historial cronológico filtrado.}
{\textbf{Pre:} existen registros previos. \textbf{Post:} se muestra el historial conforme a los filtros aplicados.}
{Should}
{La persona usuaria puede recuperar, filtrando por lote y rango de fechas, el historial completo de actividades y diagnósticos registrados.}

\subsection{Requisitos no funcionales (ISO/IEC 25010:2023)}

Los diecinueve requisitos no funcionales se organizan según las nueve características de calidad del
producto de ISO/IEC 25010:2023 \cite{iso25010}. Cada uno se enuncia con métrica, valor objetivo y
criterio de verificación observable.

{\setlength{\tabcolsep}{4pt}
\begin{longtable}{|C{1.2cm}|L{2.4cm}|L{5.8cm}|L{3.1cm}|L{2.2cm}|}
\hline
\rowcolor{grisclaro}\textbf{ID} & \textbf{Característica} & \textbf{Requisito cuantificado} & \textbf{Criterio de verificación} & \textbf{Origen} \\
\hline
\endhead
\id{RNF-01} & Adecuación funcional & La clasificación de madurez del racimo alcanza una exactitud $\geq 85\%$ y una exhaustividad $\geq 80\%$ sobre la clase ``verde'' en el conjunto de prueba. \textbf{Este umbral genérico queda precisado y sustituido por \id{RNF-IA-08}} (§9), que fija la métrica y el procedimiento de medición definitivos para el mismo componente; se conserva aquí como requisito de origen. & Matriz de confusión sobre conjunto etiquetado independiente de $\geq 200$ imágenes. \textbf{El conjunto no existe} (\id{L-05}); su construcción se incorpora al alcance de la versión como tarea previa a la verificación: recolección en campo durante dos ciclos de cosecha y etiquetado por la persona asesora técnica. Hasta su disponibilidad, el valor objetivo permanece \emph{no verificado}. & \footnotesize \id{EV-04} \\ \hline
\id{RNF-02} & Adecuación funcional & La detección de plagas alcanza una exactitud $\geq 80\%$ sobre las cinco afectaciones más frecuentes del cultivo. \textbf{Este umbral genérico queda precisado y sustituido por \id{RNF-IA-01}} (§9), que emplea macro-F1 en vez de exactitud por el desbalance de clases y fija el procedimiento de medición definitivo; se conserva aquí como requisito de origen. & Evaluación sobre conjunto etiquetado por la persona asesora técnica, sujeta a la misma condición de construcción previa declarada en \id{RNF-01} y en \id{L-05}. Hasta entonces, valor objetivo \emph{no verificado}. & \footnotesize \id{EV-01}, \id{EV-02} \\ \hline
\id{RNF-03} & Eficiencia de desempeño & Las operaciones de registro y consulta responden en $\leq 3$ s para el percentil 95, con hasta 20 sesiones concurrentes. El valor se deriva de la población declarada de 62 personas (\id{EV-12}), de las cuales el 25,8\% dispone de señal permanente en el lote: 16 sesiones simultáneas como cota superior realista, redondeada a 20 para absorber picos en el campamento. & Prueba de carga con 20 usuarios virtuales. & \footnotesize \id{EV-12} \\ \hline
\id{RNF-04} & Eficiencia de desempeño & El diagnóstico por imagen entrega resultado en $\leq 10$ s con ancho de banda $\geq 5$ Mbps. & Tiempo medio sobre 20 análisis consecutivos. & \footnotesize \id{EV-01}, \id{EV-02} \\ \hline
\id{RNF-05} & Eficiencia de desempeño & El registro GPS captura la ubicación con precisión $\leq 10$ m y período de muestreo $\leq 30$ s. & Contraste de 10 puntos con referencia geodésica. & \footnotesize \id{EV-03} \\ \hline
\id{RNF-06} & Compatibilidad & La exportación hacia la planta extractora se emite en formato CSV o JSON conforme al esquema documentado, sin campos de identificación personal. & Validación del archivo exportado contra el esquema publicado. & \footnotesize \id{EV-04} \\ \hline
\id{RNF-07} & Compatibilidad & La aplicación móvil opera en Android 9.0 o superior; el panel web, en las dos últimas versiones estables de Chrome, Edge y Firefox. & Instalación y ejecución sin error en los entornos declarados. & \footnotesize Derivado \\ \hline
\id{RNF-08} & Interacción de usuario & Una persona sin experiencia previa completa el registro de labor y el análisis de imagen con una tasa de éxito $\geq 90\%$ tras una inducción $\leq 10$ minutos. & Prueba de usabilidad con $\geq 5$ participantes de perfil trabajador agrícola. & \id{EV-05}, \id{EV-06}, \id{EV-07} \\ \hline
\id{RNF-09} & Interacción de usuario & La totalidad de la interfaz se presenta en español, con etiquetas correspondientes a la terminología del glosario y sin abreviaturas técnicas no definidas. & Inspección de la totalidad de las pantallas contra el glosario de la §1.3. & \id{EV-01}, \id{EV-07} \\ \hline
\id{RNF-10} & Fiabilidad & El servicio de respaldo presenta una disponibilidad mensual $\geq 99\%$, equivalente a una indisponibilidad $\leq 7$ h 12 min (720 h $\times$ 1\%). & Monitoreo de disponibilidad durante 30 días. & \footnotesize Derivado \\ \hline
\id{RNF-11} & Fiabilidad & La aplicación opera sin conexión y sincroniza el 100\% de los registros pendientes al restablecerse la conectividad, sin pérdida ni duplicación. & 20 operaciones ejecutadas sin conexión, verificadas tras la sincronización. & \id{EV-07}, \id{EV-10} \\ \hline
\id{RNF-12} & Fiabilidad & Ante pérdida de señal GPS, el sistema conserva el registro parcial y lo marca como de precisión degradada, sin descartar la jornada. & Simulación de pérdida de señal durante una jornada de polinización. & \footnotesize \id{EV-03} \\ \hline
\id{RNF-13} & Seguridad & Las credenciales se almacenan mediante función de derivación de clave con sal; la totalidad del tráfico emplea TLS 1.2 o superior. & Inspección del almacén de credenciales y análisis del tráfico. & \footnotesize Derivado \\ \hline
\id{RNF-14} & Seguridad & Los datos de geolocalización se cifran en reposo y su acceso queda registrado en bitácora con identificación de la persona usuaria, fecha y motivo. & Auditoría de la bitácora tras 10 accesos de prueba. & \footnotesize \cite{lopdp} Art. 37 \\ \hline
\id{RNF-15} & Mantenibilidad & El sistema presenta arquitectura modular con cobertura de pruebas automatizadas $\geq 60\%$ en la capa de lógica de negocio. & Reporte de cobertura de la herramienta de integración continua. & \footnotesize Derivado \\ \hline
\rowcolor{grisclaro}\id{RNF-16} & Interacción de usuario & \textbf{Explicabilidad} --- véase la especificación ampliada en la §3.3.1. & Prueba con 6 participantes (3 técnicos, 3 no técnicos). & \footnotesize \id{EV-01}, \id{EV-02}, \id{EV-07} \\ \hline
\id{RNF-17} & Flexibilidad \newline \footnotesize(reemplazabilidad) & La migración del servicio de inferencia a un proveedor alternativo no requiere modificar el cliente móvil, al mediar una interfaz de servicio desacoplada. & Sustitución del proveedor en entorno de prueba sin cambios en el cliente. & \footnotesize Derivado \\ \hline
\id{RNF-18} & Flexibilidad \newline \footnotesize(adaptabilidad) & Los umbrales de madurez, de porcentaje de fruta verde y de ventana corte-entrega son parametrizables por variedad y por lote sin recompilar la aplicación. & Modificación de los tres umbrales desde la configuración y verificación de su efecto. & \footnotesize \id{EV-04} \\ \hline
\id{RNF-19} & Seguridad física \newline \footnotesize(\emph{safety}) & Toda recomendación de control fitosanitario emitida por el sistema (\id{RF-09}) incluye, de forma no omisible, el equipo de protección personal exigido para el producto recomendado y el intervalo de reingreso al lote tras la aplicación. El sistema no presenta la recomendación si alguno de los dos datos falta en la ficha del producto. & Revisión de las recomendaciones generadas para el catálogo completo de productos: ninguna se presenta sin ambos datos. & \footnotesize \id{EV-02} \\ \hline
\caption{Requisitos no funcionales cuantificados según ISO/IEC 25010:2023}
\end{longtable}}

\subsubsection{RNF-16 --- Requisito de explicabilidad de los componentes de inteligencia artificial}

Conforme al gatekeeper G7 de la guía entonces aplicable a la Entrega 3 (2A), todo componente basado en inteligencia
artificial recibe tratamiento de explicabilidad como requisito no funcional, siguiendo la taxonomía
de Chazette y Schneider \cite{chazette2020} y el marco de calidad de Chazette, Brunotte y Speith
\cite{chazette2021,chazette2022}.

\textbf{Justificación empírica.} La necesidad de explicación no es una imposición normativa sino un
hallazgo de campo. La persona administradora general señaló el valor del sistema para
\emph{``las personas que no tienen tanto conocimiento en el tema de plagas''} (\id{EV-01}); el
asistente de administración anticipó que la dificultad de adopción se concentraría en el personal de
labores manuales (\id{EV-07}); y en la tercera ronda una persona estudiante señaló que necesitaría
\emph{``Primero, una imagen donde me compruebe sobre ese problema''} (\id{ENTR-11}), otra que
\emph{``No tendría mucha confianza, pero si con el tiempo veo que la información que me da es correcta,
volvería mi confianza.''} (\id{ENTR-10}), y una persona profesional del área agrícola dijo esperar que
\emph{``sea precisa en qué es lo que necesita la planta''} (\id{ENTR-14}). Estas voces coinciden en que
el valor del diagnóstico automático depende de que la persona destinataria pueda comprobarlo y actuar
sobre él.

\begin{longtable}{|L{3.4cm}|L{12.0cm}|}
\hline
\rowcolor{grisclaro}\multicolumn{2}{|l|}{\textbf{RNF-16 --- Explicabilidad del diagnóstico asistido por IA}}\\
\hline
\textbf{Componentes alcanzados} & \id{RF-07} (detección de plagas), \id{RF-08} (diagnóstico nutricional) y \id{RF-21} (clasificación de madurez). \\ \hline
\textbf{Requisito} & Para cada decisión automática, el sistema debe presentar una explicación que incluya: (i) el rasgo visual que sustentó la clasificación, expresado en la terminología del dominio; (ii) el nivel de confianza en escala cualitativa de tres niveles; y (iii) la acción recomendada. La explicación se presenta en $\leq 2$ s tras el resultado, con una extensión $\leq 60$ palabras, en español y sin terminología no incluida en el glosario. \\ \hline
\textbf{Tipos de explicación} & \emph{Causal} para el diagnóstico nutricional (qué signo foliar indica qué deficiencia); \emph{por contraste} para la clasificación de madurez (por qué maduro y no verde, con referencia al umbral de la variedad); \emph{por ejemplos} para la detección de plagas (imagen de referencia de la afectación identificada). \\ \hline
\textbf{Métrica} & Comprensión reportada $\geq 4/5$ en escala Likert por participantes de perfil no técnico, y tiempo hasta la decisión de acción $\leq 60$ s. \\ \hline
\textbf{Criterio de verificación} & Prueba de validación con 6 participantes por ronda (3 de perfil técnico y 3 de perfil no técnico), replicando el tamaño muestral de \cite{chazette2022}; se mide la cobertura de dimensiones del marco, la satisfacción percibida y el acuerdo entre rondas. \\ \hline
\textbf{Origen} & \footnotesize \id{EV-01}, \id{EV-07}, \id{ENTR-10}, \id{ENTR-11}, \id{ENTR-14} \\ \hline
\caption{RNF-16 --- Especificación ampliada del requisito de explicabilidad}
\end{longtable}

\subsection{Requisitos legales}

El SIMPA trata datos personales de personas trabajadoras en relación de dependencia, incluyendo
datos de geolocalización obtenidos de forma continua durante la jornada. Esta condición hace que el
cumplimiento de la LOPDP \cite{lopdp} no sea accesorio sino estructural al diseño del sistema.

\begin{longtable}{|C{1.5cm}|C{1.8cm}|L{6.0cm}|L{5.4cm}|}
\hline
\rowcolor{grisclaro}\textbf{ID} & \textbf{Artículo} & \textbf{Obligación} & \textbf{Requisito que la implementa} \\
\hline
\endhead
\id{RL-01} & Art. 7 & Consentimiento libre, específico e informado para el tratamiento. & Aceptación explícita y registrada en el alta de la persona usuaria, con constancia de fecha y versión del texto aceptado. \\ \hline
\id{RL-02} & Art. 10 & Principio de minimización: solo los datos necesarios para la finalidad declarada. & \id{RF-32} limita la exportación al conteo de racimos y peso promedio, excluyendo datos personales. \\ \hline
\id{RL-03} & Art. 7, 10 & Consentimiento específico para el tratamiento de datos de geolocalización. & \id{RF-15} condiciona la visualización de recorridos a la existencia de consentimiento vigente y revocable. \\ \hline
\id{RL-04} & Art. 12 & Derecho de acceso de la persona titular a sus datos. & Implementado por \id{RF-40} (exportación de los datos personales por la propia titular). \id{RF-29} cubre además la consulta del avance. \\ \hline
\id{RL-05} & Art. 13 & Derecho de rectificación. & Implementado por \id{RF-41} (corrección con conservación del valor anterior en bitácora). \\ \hline
\id{RL-06} & Art. 14 & Derecho de eliminación. & Implementado por \id{RF-42} (supresión con disociación del histórico productivo). \\ \hline
\id{RL-07} & Art. 37 & Medidas de seguridad apropiadas. & \id{RNF-13} y \id{RNF-14}: cifrado en tránsito y en reposo, control de acceso por rol. \\ \hline
\id{RL-08} & Art. 40 & Registro de las actividades de tratamiento. & \id{RNF-14}: bitácora de acceso a datos de geolocalización con identificación, fecha y motivo. \\ \hline
\caption{Trazabilidad de requisitos legales: LOPDP $\rightarrow$ requisito del sistema}
\end{longtable}

\vspace{0.3cm}
\noindent\fbox{\parbox{0.97\textwidth}{%
\small\textbf{Nota sobre la aceptación del rastreo por parte del personal.} La aplicación ampliada
del cuestionario ($n=62$, \id{EV-12}) registró que el 74{,}2\% de las personas respondientes no
manifiesta preocupación por el registro de su recorrido. Sin embargo, \textbf{el 25{,}8\% restante
sí expresa alguna preocupación} (17{,}7\% ``me preocupa un poco'' y 8{,}1\% ``me preocupa mucho'').
Este resultado no autoriza a tratar el rastreo como una funcionalidad pacíficamente aceptada: una de
cada cuatro personas destinatarias manifiesta reservas. En consecuencia, \id{RL-03} exige
consentimiento específico, \textbf{revocable en cualquier momento y sin consecuencia laboral}, y
\id{RF-15} condiciona la visualización de recorridos a su vigencia. La revocación no impide el uso
del resto del sistema.}}

\subsection{Restricciones de diseño}

\begin{longtable}{|C{1.5cm}|L{6.8cm}|L{6.6cm}|}
\hline
\rowcolor{grisclaro}\textbf{ID} & \textbf{Restricción} & \textbf{Justificación / Origen} \\
\hline
\endhead
\id{RD-01} & El cliente de campo se ejecuta sobre dispositivos Android 9.0 o superior. & El 88,7\% del personal dispone de teléfono inteligente propio (\id{EV-12}), lo que hace de Android la plataforma de mayor cobertura sin costo de equipamiento. La restricción \textbf{no supone} que la totalidad del personal disponga de uno: el 11,3\% restante queda cubierto por \id{RD-10} y \id{RF-35} (registro delegado). Véase la §2.7. \\ \hline
\id{RD-02} & La aplicación opera sin conexión y sincroniza de forma diferida. & La conexión estable se limita al campamento (\id{EV-07}) y el 74,2\% del personal declara no disponer de señal permanente en el lote (\id{EV-12}, Tabla 7). El 25,8\% que sí la declara no basta para asumir conectividad continua como precondición de operación. \\ \hline
\id{RD-03} & El sistema depende de la cámara y del receptor GPS del dispositivo. & El análisis por imagen y el conteo georreferenciado son funciones nucleares. \\ \hline
\id{RD-04} & Los datos se emplean exclusivamente con fines académicos, bajo consentimiento firmado y con anonimización previa a toda publicación. & Marco ético del proyecto y LOPDP \cite{lopdp}. \\ \hline
\id{RD-05} & El respaldo reside en la nube con API REST; la inferencia de IA y la consolidación requieren conectividad. & Decisión de arquitectura: procesamiento centralizado del modelo. \\ \hline
\id{RD-06} & La interfaz debe ser operable por personas con baja alfabetización tecnológica, en español y sin dependencia de lectura extensa. & Perfil del personal de labores manuales (\id{EV-05}, \id{EV-06}, \id{EV-07}). \\ \hline
\id{RD-07} & La clasificación de madurez no puede fundarse en el color del fruto. & Criterio técnico expreso de la planta extractora: el color induce a error (\id{EV-08}). \\ \hline
\id{RD-08} & Los umbrales de decisión deben ser parametrizables por variedad. & Los umbrales de madurez difieren entre \emph{guineensis} e híbrido (\id{EV-04}). \\ \hline
\id{RD-09} & Ninguna exportación hacia terceros puede incluir datos personales ni información económica de la organización. & Principio de minimización (\id{RL-02}) y criterio expreso de la contraparte (\id{EV-04}). \\ \hline
\id{RD-10} & El sistema debe admitir el registro delegado del avance para el personal que no dispone de dispositivo propio, sin que ello degrade su consolidado ni su remuneración por avance. & El 11{,}3\% de las personas respondientes declara no usar teléfono inteligente (\id{EV-12}); excluirlas comprometería la equidad del registro. \\ \hline
\caption{Restricciones de diseño del SIMPA}
\end{longtable}

% =====================================================================
%  4. MODELADO DEL SISTEMA CON UML
% =====================================================================
\input{seccion4_uml}

% =====================================================================
%  5. PRIORIZACION Y TRAZABILIDAD EXTENDIDA
% =====================================================================
\input{seccion5_priorizacion}

% =====================================================================
%  6. MVP  y  7. CONCLUSIONES
% =====================================================================
\input{seccion6_7_mvp_conclusiones}

% =====================================================================
%  8. DIAGRAMAS DE FLUJO DE DATOS  (incorporados en la PE5)
% =====================================================================
\input{seccion8_dfd}

% =====================================================================
%  9. COMPONENTES DE INTELIGENCIA ARTIFICIAL  (incorporado en la PE5)
% =====================================================================
\input{seccion9_ia}

% =====================================================================
%  BIBLIOGRAFIA
% =====================================================================
\newpage
% Las siguientes referencias sustentan el protocolo experimental y el
% plan de analisis estadistico desarrollados en las secciones 5 y 6.
\nocite{wohlin2012,runeson2009,molleri2020,hennink2017,cohen1960,wilkinson2016,nosek2018}

% Estilo IEEE si esta disponible (paquete texlive-publishers);
% El estilo IEEEtran se distribuye localmente para garantizar una compilación reproducible.
% en orden de aparicion y esta presente en toda instalacion base.
\bibliographystyle{IEEEtran}
\bibliography{referencias}

% =====================================================================
%  APENDICES
% =====================================================================
% =====================================================================
%  APENDICES A -- E
% =====================================================================

\newpage
\appendix
\section{Plan del proyecto de Ingeniería de Requisitos}

\subsection{Identificación de las partes interesadas}

\begin{longtable}{|C{1.6cm}|L{4.4cm}|C{1.8cm}|C{1.6cm}|C{1.6cm}|C{2.6cm}|}
\hline
\rowcolor{grisclaro}\textbf{Código} & \textbf{Rol} & \textbf{Organización} & \textbf{Tipo} & \textbf{Influencia} & \textbf{Evidencia} \\
\hline
\endhead
\id{ENTR-01} & Administrador general & Palmicultora M & Primario & Alta & \id{EV-01} \\ \hline
\id{ENTR-02} & Asesor técnico & Palmicultora M & Primario & Alta & \id{EV-02} \\ \hline
\id{ENTR-03} & Jefe de polinización & Palmicultora M & Primario & Alta & \id{EV-03} \\ \hline
\id{ENTR-04} & Supervisor de técnicos agrícolas & Extractora R & Externo & Alta & \id{EV-04} \\ \hline
\id{ENTR-05} & Trabajador agrícola & Palmicultora M & Operativo & Baja & \id{EV-05} \\ \hline
\id{ENTR-06} & Trabajador agrícola & Palmicultora M & Operativo & Baja & \id{EV-06} \\ \hline
\id{ENTR-07} & Asistente de administración & Palmicultora M & Primario & Media & \id{EV-07} \\ \hline
\id{ENTR-08} & Técnico agrícola & Extractora R & Externo & Media & \id{EV-08} \\ \hline
\caption{Identificación de las partes interesadas entrevistadas}
\end{longtable}

\subsection{Cronograma de actividades}

\begin{longtable}{|C{2.0cm}|L{6.6cm}|L{6.4cm}|}
\hline
\rowcolor{grisclaro}\textbf{Semana} & \textbf{Actividad} & \textbf{Producto} \\
\hline
\endhead
1--3 & Identificación del problema y del sistema objetivo & Propuesta de proyecto \\ \hline
3--4 & Primera ronda de elicitación (\id{EV-01} a \id{EV-03}) y observación de campo & Entrega 1 (1A) \\ \hline
5--6 & Elaboración del paquete ético y obtención del aval institucional & Documentos A.1--A.13 y Categoría C \\ \hline
6--9 & Modelado UML inicial y formalización de requisitos & Diagramas y plantillas de RF \\ \hline
9--10 & Consolidación del ERS parcial & Entrega 2 (1B) \\ \hline
11 & Adenda ética y registro del protocolo experimental en OSF & Adenda y comprobante OSF \\ \hline
12 & Segunda ronda de campo (\id{EV-04} a \id{EV-08}) y ampliación del cuestionario & Grabaciones, transcripciones, $n=62$ \\ \hline
13 & Consolidación del ERS completo y del MVP & \textbf{Entrega 3 (2A)} \\ \hline
14--16 & Ejecución del componente empírico y análisis & Resultados y borrador de manuscrito \\ \hline
17 & Defensa final & Entrega 4 (2B) \\ \hline
\caption{Cronograma de actividades de ingeniería de requisitos}
\end{longtable}

\subsection{Técnicas de elicitación y justificación}

\begin{longtable}{|L{3.6cm}|C{1.6cm}|L{9.4cm}|}
\hline
\rowcolor{grisclaro}\textbf{Técnica} & \textbf{Sesiones} & \textbf{Justificación} \\
\hline
\endhead
Entrevista semiestructurada & 8 & Permite profundizar en el conocimiento experto y adaptarse a respuestas inesperadas. Es la técnica indicada cuando el dominio contiene conocimiento tácito \cite{ferrari2016}, situación confirmada en el criterio de madurez del racimo. \\ \hline
Cuestionario estructurado & 2 aplicaciones & Permite cuantificar la utilidad percibida y las condiciones de operación en una población que no puede ser entrevistada en su totalidad. Instrumento idéntico entre rondas para preservar la comparabilidad. \\ \hline
Observación directa & 1 & Documenta el proceso tal como ocurre, sin la mediación del relato. Registrada en \id{EV-09}. \\ \hline
Análisis de documentos & 2 tipos & Los planes semanales (\id{EV-11}) revelaron la estructura de planificación y las unidades reales de medida del trabajo. La hoja de liquidación (\id{EV-13}) hizo explícito el catálogo de tarifas por labor, que ninguno de los ocho participantes enunció pese a operar bajo él. Es la técnica que más requisitos aportó por unidad de esfuerzo. \\ \hline
Prototipado & 8 pantallas & Los prototipos de alta fidelidad (\id{MU-01} a \id{MU-08}) permiten validar la interfaz con personal de baja alfabetización tecnológica antes de construir. \\ \hline
\caption{Técnicas de elicitación aplicadas y su justificación}
\end{longtable}

\subsection{Declaración de conflicto de interés}

Se declara que dos de las personas participantes mantienen relación de consanguinidad con el
analista líder del equipo: el participante \id{ENTR-02} en primer grado y el participante
\id{ENTR-07} en segundo grado. Ambas entrevistas fueron conducidas por el propio analista líder.

Se aplicaron seis medidas de mitigación:

\begin{enumerate}[nosep]
  \item Presencia de un segundo integrante del equipo en calidad de registrador durante ambas sesiones.
  \item Conservación íntegra de las grabaciones y transcripciones, disponibles para auditoría docente.
  \item Uso estricto del banco de preguntas aprobado, sin preguntas improvisadas fuera del guion.
  \item Verificación de los requisitos derivados por integrantes con rol de verificador que no participaron en esas sesiones.
  \item Triangulación obligatoria de todo hallazgo con al menos una de las otras seis entrevistas.
  \item \textbf{Ningún requisito del presente ERS se sustenta de forma exclusiva en estas dos fuentes.}
\end{enumerate}

\noindent
La sexta medida es verificable en la matriz de trazabilidad: toda fila que referencia \id{EV-02} o
\id{EV-07} referencia además al menos una evidencia adicional independiente.

% =====================================================================
\newpage
\section{Trabajo de campo y evidencias}

\subsection{Catálogo de evidencias}

\footnotesize
\begin{longtable}{|C{1.2cm}|L{2.6cm}|C{1.7cm}|C{1.8cm}|L{4.2cm}|C{1.3cm}|C{1.1cm}|}
\hline
\rowcolor{grisclaro}\textbf{ID-EV} & \textbf{Tipo} & \textbf{Fecha} & \textbf{Participante} & \textbf{Requisitos derivados} & \textbf{Zona} & \textbf{Cons.} \\
\hline
\endhead
\id{EV-01} & Entrevista (audio) & 23/05/2026 & \id{ENTR-01} & \id{RF-01} a \id{RF-06}, \id{RF-09}, \id{RF-12}, \id{RF-15} a \id{RF-20} & [R] & Sí \\ \hline
\id{EV-02} & Entrevista (video) & 23/05/2026 & \id{ENTR-02} & \id{RF-02}, \id{RF-05}, \id{RF-08}, \id{RF-11}, \id{RF-17}, \id{RF-18} & [R] & Sí \\ \hline
\id{EV-03} & Entrevista (video) & 23/05/2026 & \id{ENTR-03} & \id{RF-13}, \id{RF-14}, \id{RF-15}, \id{RNF-05} & [R] & Sí \\ \hline
\id{EV-04} & Entrevista (video) & 28/07/2026 & \id{ENTR-04} & \id{RF-10}, \id{RF-21} a \id{RF-25}, \id{RF-32}, \id{RF-33}, \id{RNF-18} & [R] & Sí \\ \hline
\id{EV-05} & Entrevista (video) & 28/07/2026 & \id{ENTR-05} & \id{RF-04}, \id{RF-28}, \id{RF-35} & [R] & Sí \\ \hline
\id{EV-06} & Entrevista (video) & 28/07/2026 & \id{ENTR-06} & \id{RF-04}, \id{RF-28}, \id{RF-29} & [R] & Sí \\ \hline
\id{EV-07} & Entrevista (video) & 02/08/2026 & \id{ENTR-07} & \id{RF-29} a \id{RF-31}, \id{RNF-11} & [R] & Sí \\ \hline
\id{EV-08} & Entrevista (video) & 28/07/2026 & \id{ENTR-08} & \id{RF-21}, \id{RF-22}, \id{RF-24}, \id{RD-07} & [R] & Sí \\ \hline
\id{EV-09} & Observación de campo & 23/05/2026 & Equipo AHMRV & \id{RF-01}, \id{RF-02}, \id{RF-05} & [P] & N/A \\ \hline
\id{EV-10} & Cuestionario piloto & 30/06/2026 & $n=4$ & Validación de \id{RF-04}, \id{RF-07}, \id{RF-12} & [P] & Anónimo \\ \hline
\id{EV-11} & Documentos de la organización & 29/06 a 24/07/2026 & Palmicultora M & \id{RF-02}, \id{RF-03}, \id{RF-19}, \id{RF-26} a \id{RF-28}, \id{RF-34} & [P] + [R] & N/A \\ \hline
\id{EV-12} & Cuestionario ampliado & 30/06/2026 & $n=62$ & \id{RF-12}, \id{RF-15}, \id{RF-35}, \id{RD-10}, \id{RNF-08} & [P] & Anónimo \\ \hline
\id{EV-13} & Hoja de liquidación semanal & 03--07/08/2026 & Palmicultora M & \id{RF-36} a \id{RF-39} & [P] + [R] & N/A \\ \hline
\caption{Catálogo de evidencias de campo. [P] = zona pública, [R] = zona restringida cifrada}
\label{tab:catalogo}
\end{longtable}

\noindent
\textbf{Nota sobre la renumeración.} Las evidencias \id{EV-09} y \id{EV-10} correspondían a
\id{EV-04} y \id{EV-05} en la Entrega 2 (1B). El desplazamiento permitió asignar identificadores
contiguos a las entrevistas de la segunda ronda. La equivalencia se registra en \id{CHANGELOG.md}.

\subsection{Documentos de la organización}

Se obtuvieron dos tipos documentales distintos, ambos correspondientes a la operación real de la
organización cliente durante el período del proyecto.

\begin{longtable}{|C{1.4cm}|L{4.2cm}|C{2.4cm}|L{6.4cm}|}
\hline
\rowcolor{grisclaro}\textbf{ID-EV} & \textbf{Tipo documental} & \textbf{Período} & \textbf{Aporte a la especificación} \\
\hline
\endhead
\id{EV-11} & Plan semanal de labores (presupuesto por cumplir y cumplido) & Semanas 27 a 30 de 2026 & Estructura de lotes, catálogo de labores, unidades de avance, patrón de arrastre de labor incompleta, requerimientos de insumos \\ \hline
\id{EV-13} & Hoja de liquidación de presupuesto contra ejecución & Semana 32 de 2026 & Código de labor, unidad de medida, tarifa vigente, distinción entre cantidad facturada y ejecutada, grupo de labor no presupuestada, flete como costo independiente, aprobación de la administración \\ \hline
\caption{Documentos de la organización incorporados como evidencia}
\end{longtable}

\noindent
\textbf{Relación entre ambos documentos.} \id{EV-11} declara metas; \id{EV-13} liquida su ejecución.
El par constituye el ciclo completo de planificación y control que el sistema debe reproducir, y
sustenta los requisitos \id{RF-26}, \id{RF-27} y \id{RF-37} a \id{RF-39}. Su obtención en momentos
distintos del proyecto explica que los requisitos de liquidación aparezcan al final de la numeración
pese a corresponder a un proceso central de la organización.

\noindent
\textbf{Tratamiento aplicado.} Conforme a la §8.4 de la guía, la copia depositada en la zona pública
tiene tachados de forma irreversible los valores de tarifa, los importes por labor y los totales de
gasto, así como el nombre propio que figuraba en la columna de observaciones, sustituido por el
código de participante correspondiente. Los originales íntegros residen en el contenedor cifrado.

\subsection{Instrumentos aplicados}

Los instrumentos residen en \id{06\_Experimento/instrumentos/}. Se aplicaron cuatro guiones
diferenciados por perfil, decisión metodológica que permite ajustar el vocabulario y la profundidad
técnica al interlocutor:

\begin{longtable}{|C{2.0cm}|L{5.0cm}|C{2.0cm}|L{5.4cm}|}
\hline
\rowcolor{grisclaro}\textbf{Guion} & \textbf{Perfil destinatario} & \textbf{Aplicado a} & \textbf{Foco temático} \\
\hline
\endhead
\id{G-01} & Dirección y asesoría técnica & \id{EV-01}, \id{EV-02} & Manejo del cultivo, nutrición, sanidad, estimación \\ \hline
\id{G-02} & Jefatura de polinización & \id{EV-03} & Antesis, calidad de la labor, control del personal \\ \hline
\id{G-03} & Personal de la extractora & \id{EV-04}, \id{EV-08} & Criterios de calidad, calificación, penalización \\ \hline
\id{G-04} & Personal de campo y asistencia & \id{EV-05} a \id{EV-07} & Jornada, registro, reporte de afectaciones, tecnología \\ \hline
\id{C-01} & Cuestionario (idéntico en ambas rondas) & \id{EV-10}, \id{EV-12} & Perfil, registro actual, dificultades, utilidad percibida, aceptación \\ \hline
\caption{Instrumentos de recolección aplicados}
\end{longtable}

\subsection{Codificación temática y saturación}

La codificación temática de las ocho transcripciones reside en
\id{07\_Datos/datos\_crudos/}\allowbreak\id{codificacion.csv}, con las columnas Fragmento, Código,
Categoría, Requisito derivado, ID de evidencia y Analista codificador.

\begin{longtable}{|C{1.6cm}|C{2.4cm}|C{2.4cm}|C{2.4cm}|L{5.0cm}|}
\hline
\rowcolor{grisclaro}\textbf{Entrevista} & \textbf{Códigos totales} & \textbf{Códigos nuevos} & \textbf{Acumulado} & \textbf{Observación} \\
\hline
\endhead
\id{EV-01} & 24 & 24 & 24 & Primera fuente: todo es nuevo \\ \hline
\id{EV-02} & 21 & 10 & 34 & Alta redundancia con \id{EV-01}: 52\% de los fragmentos reiteran códigos previos \\ \hline
\id{EV-03} & 14 & 6 & 40 & Aporta el dominio de la polinización asistida \\ \hline
\id{EV-04} & 26 & \textbf{17} & 57 & \textbf{Repunte}: dominio nuevo (calidad de la fruta en la extractora) \\ \hline
\id{EV-05} & 9 & 4 & 61 & Confirma el registro manual en papel \\ \hline
\id{EV-06} & 11 & 1 & 62 & \textbf{Mínimo de la serie}: 9\% de aportación \\ \hline
\id{EV-07} & 15 & 4 & 66 & Aporta la latencia de verificación y el canal informal \\ \hline
\id{EV-08} & 18 & 2 & 68 & Triangula \id{EV-04}; 11\% de aportación \\ \hline
\caption{Aportación de códigos nuevos por entrevista (curva de saturación)}
\label{tab:saturacion}
\end{longtable}

\begin{figure}[htbp]
\centering
\includegraphics[width=0.96\textwidth]{img/curva_saturacion}
\caption{Curva de saturación temática: aportación de códigos nuevos por entrevista, en orden cronológico de recolección.}
\label{fig:saturacion}
\end{figure}

\noindent
\textbf{Lectura de la curva.} La aportación de códigos nuevos desciende de forma sostenida hasta
\id{EV-03}, con \id{EV-02} reiterando el 52\% de sus fragmentos y \id{EV-03} el 57\%. Ese descenso
sugería saturación al cierre de la primera ronda.

El repunte de \id{EV-04} (17 códigos nuevos sobre 26 fragmentos, un 65\% de aportación) demuestra
que dicha saturación era \textbf{aparente}: correspondía al agotamiento de un dominio ---el manejo
del cultivo--- y no del problema completo. La incorporación de la organización externa abrió un
dominio no explorado: los criterios de calidad con que se valora económicamente la producción.

La caída posterior es igual de informativa. \id{EV-06} aporta un solo código nuevo sobre once
fragmentos (9\%) y \id{EV-08} aporta dos sobre dieciocho (11\%), pese a pertenecer este último al
mismo dominio recién abierto. Es decir, el dominio de la calidad de la fruta \emph{sí} alcanzó
saturación, pero con dos fuentes, no con ocho entrevistas en total.

Este comportamiento ---saturación por dominio y no por número absoluto de entrevistas--- es
coherente con la distinción entre saturación de código y saturación de significado de Hennink
\emph{et al.} \cite{hennink2017}. Se reportará como hallazgo metodológico en el manuscrito, con la
implicación práctica de que el criterio de suficiencia muestral debe evaluarse por dominio de
conocimiento y no sobre el agregado.

\subsection{Zonas de evidencia y verificación}

Conforme a la §4.1 de la guía, la evidencia se organiza en dos zonas. La zona pública \textbf{[P]}
contiene transcripciones anonimizadas, copias enmascaradas de consentimientos, fotografías sin
rostros ni coordenadas GPS, respuestas del cuestionario sin columnas identificativas y las fichas
técnicas de cada archivo multimedia. La zona restringida \textbf{[R]} contiene los originales
identificables, cifrados con AES-256 en
\id{02\_Evidencias/00\_Restringido/}\allowbreak\id{evidencias\_restringidas\_2A.7z}.

El hash SHA-256 de cada archivo se calculó \emph{antes} del cifrado y consta en
\id{checksums.sha256} junto con la ficha técnica pública \id{fichas\_tecnicas.csv}, que registra
nombre de archivo, tipo, fecha, código de participante, duración, códec, tamaño y hash. La
contraseña del contenedor se entrega únicamente al docente por el espacio de la actividad en el
sistema de gestión académica.

\noindent\fbox{\parbox{0.97\textwidth}{%
\small\textbf{Nota técnica sobre Git LFS.} Los archivos multimedia y el contenedor cifrado se
versionan mediante Git Large File Storage. Una clonación sin \id{git lfs pull} muestra archivos de
133 bytes correspondientes a los punteros, no al contenido. Las instrucciones de recuperación
constan en el \id{README.md} de la raíz del repositorio.}}

% =====================================================================
\newpage
\section{Clasificación y priorización completa}

La tabla completa de clasificación de los 42 requisitos funcionales, con su prioridad MoSCoW, su
categoría de Kano, los cuatro factores del cálculo WSJF y la justificación de cada asignación,
reside en \id{04\_Trazabilidad/}\allowbreak\id{priorizacion\_moscow\_kano.csv}. La §5.3 del cuerpo del documento
desarrolla la argumentación de las categorías extremas y el contraste entre las tres técnicas.

Los requisitos no funcionales y las restricciones de diseño se clasifican como sigue.

\begin{longtable}{|C{2.0cm}|C{2.0cm}|L{5.0cm}|L{5.4cm}|}
\hline
\rowcolor{grisclaro}\textbf{ID} & \textbf{Prioridad} & \textbf{Característica / Ámbito} & \textbf{Razón} \\
\hline
\endhead
\id{RNF-01}, \id{RNF-02} & Must & Adecuación funcional & La exactitud del clasificador determina la utilidad del sistema \\ \hline
\id{RNF-03} & Should & Eficiencia de desempeño & Deseable; sin umbral crítico en la operación de campo \\ \hline
\id{RNF-04} & Must & Eficiencia de desempeño (IA) & Un diagnóstico lento no se usa en campo \\ \hline
\id{RNF-05} & Should & Eficiencia (precisión GPS) & Condiciona la confiabilidad del conteo \\ \hline
\id{RNF-06} & Should & Compatibilidad & Requerido para el intercambio con la extractora \\ \hline
\id{RNF-07} & Should & Compatibilidad & Cobertura de versiones de plataforma \\ \hline
\id{RNF-08}, \id{RNF-09} & Must & Interacción de usuario & Condicionan la adopción por personal de baja alfabetización digital \\ \hline
\id{RNF-10} & Should & Fiabilidad & Disponibilidad alta; no bloqueante en v1 \\ \hline
\id{RNF-11} & Must & Fiabilidad & El 74{,}2\% del personal carece de conectividad permanente \\ \hline
\id{RNF-12} & Should & Fiabilidad & Tolerancia a degradación de señal GPS \\ \hline
\id{RNF-13}, \id{RNF-14} & Must & Seguridad & Exigidos por los artículos 37 y 40 de la LOPDP \\ \hline
\id{RNF-19} & Must & Seguridad física & Condiciona la recomendación fitosanitaria de \id{RF-09}: su omisión expone físicamente al personal \\ \hline
\id{RNF-15} & Could & Mantenibilidad & Meta de calidad interna \\ \hline
\id{RNF-16} & Must & Explicabilidad & Obligatorio para todo componente de IA; gatekeeper G7 \\ \hline
\id{RNF-17} & Should & Portabilidad & Evita dependencia de un proveedor de inferencia \\ \hline
\id{RNF-18} & Must & Flexibilidad & Los umbrales varían por variedad; prerrequisito de \id{RF-21} \\ \hline
\id{RD-01} a \id{RD-10} & Must & Restricciones & Impuestas por el entorno, la normativa o la contraparte \\ \hline
\caption{Clasificación de los requisitos no funcionales y las restricciones}
\end{longtable}

% =====================================================================
\newpage
\section{Matriz de trazabilidad completa}

La matriz completa, con las trece columnas exigidas por la §8.8 de la guía, reside en
\id{04\_Trazabilidad/}\allowbreak\id{matriz\_trazabilidad.csv} en formato procesable. La §5.4 del cuerpo del
documento presenta la vista reducida de sus 55 filas y el diagnóstico de cobertura correspondiente.

\subsection{Diccionario de identificadores}

\begin{longtable}{|C{2.2cm}|L{5.2cm}|L{7.0cm}|}
\hline
\rowcolor{grisclaro}\textbf{Prefijo} & \textbf{Significado} & \textbf{Rango} \\
\hline
\endhead
\id{TR-XX} & Fila de la matriz de trazabilidad & \id{TR-01} a \id{TR-55} \\ \hline
\id{CB-XX} & Criterio de aceptación BDD derivado de la ficha del requisito & \id{CB-01} a \id{CB-37} \\ \hline
\id{O-X} & Objetivo del sistema & \id{O1} a \id{O6} \\ \hline
\id{EV-XX} & Identificador de evidencia & \id{EV-01} a \id{EV-13} \\ \hline
\id{ENTR-XX} & Código de participante & \id{ENTR-01} a \id{ENTR-08} \\ \hline
\id{RF-XX} & Requisito funcional & \id{RF-01} a \id{RF-42} \\ \hline
\id{RNF-XX} & Requisito no funcional & \id{RNF-01} a \id{RNF-19} \\ \hline
\id{RD-XX} & Restricción de diseño & \id{RD-01} a \id{RD-10} \\ \hline
\id{RL-XX} & Requisito legal (LOPDP) & \id{RL-01} a \id{RL-08} \\ \hline
\id{CU-XX} & Caso de uso & \id{CU-01} a \id{CU-18} \\ \hline
\id{HU-XX} & Historia de usuario & \id{HU-01} a \id{HU-19} \\ \hline
\id{CA-XX} & Criterio de aceptación & \id{CA-01} a \id{CA-19} \\ \hline
\id{MU-XX} & Prototipo de interfaz & \id{MU-01} a \id{MU-08} \\ \hline
\id{C-XXX} & Componente de la arquitectura & Véase Figura~\ref{fig:componentes} \\ \hline
\id{L-XX} & Limitación declarada & \id{L-01} a \id{L-11} \\ \hline
\id{D-XX} & Desviación del prototipo & \id{D-01} a \id{D-03} \\ \hline
\caption{Diccionario de prefijos de identificadores empleados en el documento}
\end{longtable}

% =====================================================================
\newpage
\section{Repositorio del proyecto}

\subsection{Enlace y estructura}

\begin{center}
\url{https://github.com/gleiston-guerrero/SIMPA_ISR401}
\end{center}

\noindent
Los cinco archivos raíz exigidos ---\id{README.md}, \id{LICENSE}, \id{CITATION.cff},
\id{CHANGELOG.md} y \id{checksums.sha256}--- se encuentran completos en la raíz del repositorio.

\vspace{0.3cm}
\noindent\fbox{\parbox{0.97\textwidth}{%
\small\textbf{Declaración de desviación estructural.} La §8.1 de la guía establece que la raíz del
repositorio debe reproducir el árbol del proyecto. En este repositorio el proyecto reside en la
carpeta \id{AHMRV/}, un nivel por debajo de la raíz.

\vspace{0.2cm}
La desviación es deliberada. La ruta \id{/tree/main/AHMRV} es la declarada en la portada del
documento entregado en el sistema de gestión académica, cuya actividad se encuentra cerrada y no
admite modificación por ningún medio. Trasladar el contenido a la raíz produciría un error 404 en
ese enlace, dado que GitHub no redirige rutas eliminadas, y activaría el gatekeeper G1 con nota
máxima de 2,50 sobre 10.

\vspace{0.2cm}
Ante la disyuntiva entre una desviación estructural penalizable en rúbrica y un enlace roto que
activa un gatekeeper, el equipo optó por preservar la verificabilidad del enlace y declarar la
desviación de forma expresa. La estructura interna de \id{AHMRV/} reproduce íntegramente el árbol
de la §8.1, y el \id{README.md} de la raíz orienta hacia ella.}}

\begin{longtable}{|L{4.6cm}|L{10.8cm}|}
\hline
\rowcolor{grisclaro}\textbf{Carpeta} & \textbf{Contenido} \\
\hline
\endhead
\id{01\_ERS/} & Documento en PDF, fuente \id{.tex}, \id{referencias.bib} e imágenes vectoriales \\ \hline
\id{02\_Evidencias/} & Zona restringida cifrada, consentimientos enmascarados, transcripciones, fotografías, cuestionario, documentos de la organización, actas de validación y codificación temática \\ \hline
\id{03\_Modelado/} & Trece diagramas UML e i* en fuente PlantUML, PNG a 300 dpi y SVG; ocho prototipos de interfaz \\ \hline
\id{04\_Trazabilidad/} & \id{matriz\_trazabilidad.csv} y \id{priorizacion\_moscow\_kano.csv} \\ \hline
\id{05\_MVP/} & Enlace al repositorio del prototipo y video de demostración \\ \hline
\id{06\_Experimento/} & Protocolo, comprobante OSF, instrumentos, consignas de LLM, resultados y scripts \\ \hline
\id{07\_Publicacion/} & Borrador de manuscrito, análisis de revistas y paquete de datos para Zenodo \\ \hline
\id{08\_Etica/} & Paquete ético completo, documentos de Categoría C y adenda de la segunda ronda \\ \hline
\caption{Estructura del repositorio del proyecto}
\end{longtable}

\subsection{Instrucciones de reproducción}

\textbf{Obtención del repositorio con los archivos multimedia:}

\begin{quote}
\id{git clone https://github.com/AlanNVR/}\allowbreak\id{Villafuerte\_Grupo\_AHMRV.git}\\
\id{cd Villafuerte\_Grupo\_AHMRV}\\
\id{git lfs pull}
\end{quote}

\noindent
\textbf{Compilación del presente documento} desde \id{01\_ERS/}:

\begin{quote}
\id{pdflatex ERS\_SRS\_2A\_v1.0.tex}\\
\id{bibtex ERS\_SRS\_2A\_v1.0}\\
\id{pdflatex ERS\_SRS\_2A\_v1.0.tex}\\
\id{pdflatex ERS\_SRS\_2A\_v1.0.tex}
\end{quote}

\noindent
\textbf{Regeneración de los diagramas} desde \id{03\_Modelado/Diagramas\_UML/}:

\begin{quote}
\id{plantuml -tpng -Sdpi=300 -o png *.puml}\\
\id{plantuml -tsvg -o svg *.puml}
\end{quote}

\noindent
\textbf{Verificación de integridad} desde la raíz: \id{sha256sum -c checksums.sha256}

\subsection{Contribución del equipo}

\begin{longtable}{|L{4.4cm}|L{4.0cm}|L{6.8cm}|}
\hline
\rowcolor{grisclaro}\textbf{Integrante} & \textbf{Rol} & \textbf{Contribución principal en la 2A} \\
\hline
\endhead
Villafuerte Rosero Allan Noe & Analista líder & Coordinación, segunda ronda de elicitación, paquete ético y adenda, protocolo experimental \\ \hline
Huilcapi León Denisses Fabiola & Documentadora & Redacción del ERS, migración a \LaTeX, control de versiones \\ \hline
Rizzo Vélez Edson Nagib & Verificador & Verificación de calidad de requisitos, gestión de evidencias y fichas técnicas \\ \hline
Macías Herrera Josthyn Esteban & Modelador & Diagramas UML e i*, matriz de trazabilidad \\ \hline
Arboleda Yanza Francisco Javier & Apoyo modelado & Estructura del repositorio, priorización, diagramas \\ \hline
Alcívar Vélez Anderson Adonis & Apoyo repositorio & Transcripciones, anonimización, prototipo funcional \\ \hline
\caption{Contribución de los integrantes del equipo en la Entrega 4 (2B)}
\end{longtable}

\input{apendice_bdd}



% =====================================================================
% DECLARACIÓN DE USO DE INTELIGENCIA ARTIFICIAL — Entrega Final 2B
% =====================================================================
\input{declaracion_uso_ia}

\end{document}
